%%
% BIThesis 本科毕业设计论文模板 —— 使用 XeLaTeX 编译 The BIThesis Template for Undergraduate Thesis
% This file has no copyright assigned and is placed in the Public Domain.
% Compile with: xelatex -> biber -> xelatex -> xelatex
%%

% 请勿删除下面两行注释，以免影响编译。
% !TeX program = xelatex
% !BIB program = biber


% 开启盲审格式 blindPeerReview=true (如：[type=bachelor,blindPeerReview=true])

\documentclass[type=bachelor]{bithesis}

% 此处仅列出常用的配置。全部配置用法请见「bithesis.pdf」手册。
\BITSetup{
  cover = {
    % 封面需要「北京理工大学」字样图片，如无必要请勿修改该项。
    headerImage = images/header.pdf,
    % 封面标题需要"华文细黑"，如无必要请勿修改该项。
    xiheiFont = STXIHEI.TTF,
  },
  info = {
    title = 基于多智能体的学术多模态属性图数据集自动化构建系统,
    titleEn = {Automated Construction System for Academic Multimodal Attributed Graph Dataset Based on Multi-Agent Architecture},
    school = 计算机学院,
    major = 计算机科学与技术,
    class = 0561xxxx,
    author = xxx,
    studentId = 11xxxxxxxx,
    supervisor = xxx,
    keywords = {多模态属性图；图神经网络；学术文献；多智能体系统；数据集构建},
    keywordsEn = {Multimodal Attributed Graph; Graph Neural Network; Academic Literature; Multi-Agent System; Dataset Construction},
  },
  style = {
    betterTimesNewRoman = true,
  },
  misc = {
    tabularRowSeparation = 1.25,
  },
}

\usepackage{listings}
\usepackage{booktabs}
\usepackage{multirow}
\usepackage{graphicx}
\usepackage{amsmath}
\usepackage{hyperref}
\usepackage{xcolor}
\usepackage{tikz}
\usetikzlibrary{shapes, arrows, positioning, fit, backgrounds}

% 大部分关于参考文献样式的修改，都可以通过此处的选项进行配置。
% 详情请搜索「biblatex-gb7714-2015 文档」进行阅读。
\usepackage[
  backend=biber,
  style=gb7714-2015,
  gbalign=gb7714-2015,
  gbnamefmt=lowercase,
  gbpub=false,
  doi=false,
  url=false,
  eprint=false,
  isbn=false,
]{biblatex}

% 参考文献引用文件位于 misc/ref.bib
\addbibresource{misc/ref.bib}

% 文档开始
\begin{document}

% 标题页面
\MakeCover

% 原创性声明
\begin{blindPeerReview}
  \includepdf{misc/1_originality.pdf}\newpage
\end{blindPeerReview}

% 前置内容定义
\frontmatter
% 摘要
%%
% BIThesis 本科毕业设计论文模板 —— 使用 XeLaTeX 编译 The BIThesis Template for Undergraduate Thesis
% This file has no copyright assigned and is placed in the Public Domain.
%%

% 中英文摘要章节
\begin{abstract}

现有MAG基准数据集（MAGB、OpenMAG、MM-GRAPH）均限于电商和社交领域且仅含文本与图像双模态，未覆盖学术论文领域和公式模态。针对这一空白，本文提出AMMAG——首个面向学术论文领域的多模态属性图Benchmark数据集，覆盖计算机科学、数学和物理学三个学科共45,672篇论文，包含文本、图像、公式三种模态和59,364条引用边，并通过基于多智能体的自动化系统实现端到端构建。在Benchmark实验方面，本文在六组任务上系统评估了6种GNN模型与7种模态组合。核心发现包括：（1）三模态融合较文本单模态提升+0.45pp（20次运行，p$<$0.001），达到统计显著，但不同领域模态偏好差异显著——cs.SI获益最大（+1.27pp）而math.LO几乎不受图像影响；（2）公式密度效应揭示了通用编码器在数学公式上的局限；（3）L2归一化在Few-shot中提升8-11pp，但归一化后Image仍占拼接特征58.6\%的能量，表明简单拼接不足以解决模态失衡。AMMAG为学术MAG学习研究提供了首个标准化评估平台。

\end{abstract}

% 英文摘要章节
\begin{abstractEn}

Existing MAG benchmarks (MAGB, OpenMAG, MM-GRAPH) are limited to e-commerce and social domains with only text and image modalities, failing to cover academic literature and formula modality. We propose AMMAG --- the first multimodal attributed graph benchmark for academic literature, covering 45,672 papers across three disciplines with text, image, and formula modalities along with 59,364 citation edges, constructed via an automated multi-agent system. We systematically evaluate 6 GNN models across 7 modality combinations on six task groups. Key findings include: (1) tri-modal fusion achieves a statistically significant improvement of +0.45pp over text-only (20 runs, p<0.001), with domain-specific modality preferences --- cs.SI benefits most (+1.27pp) while math.LO is barely affected by images; (2) the formula density effect reveals limitations of general-purpose encoders on mathematical formulas; (3) L2 normalization improves few-shot performance by 8-11pp, yet images still dominate 58.6\% of concatenated feature energy, indicating simple concatenation is insufficient. AMMAG provides the first standardized evaluation platform for academic MAG learning.
\end{abstractEn}


\MakeTOC

% 正文开始
\mainmatter

% 第一章 绪论
%%
% 第一章 绪论（~5页）
%%

\chapter{绪论}

\section{研究背景与意义}

近年来，图神经网络（Graph Neural Network, GNN）作为处理图结构数据的深度学习方法，在社交网络分析\cite{hamilton2017rep_survey}、推荐系统、药物发现和交通预测等领域取得了显著的成功。图数据广泛存在于现实世界中，其中节点表示实体，边表示实体间的关系。与传统欧几里得数据（如图像、文本）不同，图数据的非规则结构使得传统卷积操作难以直接应用。针对这一挑战，研究者们提出了一系列图神经网络模型。Kipf和Welling\cite{kipf2017gcn}提出了图卷积网络（Graph Convolutional Network, GCN），通过谱域近似的方法将卷积操作推广到图结构上，实现了高效的邻居信息聚合。Veli\v{c}kovi\'{c}等人\cite{velickovic2018gat}提出了图注意力网络（Graph Attention Network, GAT），引入注意力机制自适应地学习邻居节点的重要性权重。Hamilton等人\cite{hamilton2017graphsage}提出了GraphSAGE，通过采样和聚合的策略支持归纳式学习，使其能够处理大规模动态图。Xu等人\cite{xu2019gin}从理论上分析了GNN的表达能力上界，提出了图同构网络（Graph Isomorphism Network, GIN），证明其表达能力与Weisfeiler-Lehman图同构测试等价。这些模型在节点分类、图分类、链接预测等核心图学习任务上展现了强大的性能\cite{zhou2020gnn_survey, wu2021gnn_survey}。

Benchmark数据集在推动机器学习发展方面发挥着至关重要的作用。在图像识别领域，ImageNet数据集推动了深度卷积神经网络的快速发展；在自然语言处理领域，GLUE基准促进了预训练语言模型的持续进步。类似地，图机器学习领域也迫切需要高质量、大规模、标准化的Benchmark数据集来客观评估和比较不同方法的性能。Hu等人\cite{hu2020ogb}提出了开放图基准（Open Graph Benchmark, OGB），提供了涵盖节点分类、链接预测和图分类三大任务的标准化评估框架。其中，ogbn-arxiv数据集包含约169K篇arXiv论文节点和1.16M条引用边，将论文的标题和摘要编码为128维skip-gram向量作为节点特征。ogbn-mag数据集则包含约1.9M个节点和21M条边，覆盖论文、作者、机构和研究领域四种节点类型。OGB极大地促进了图学习研究的标准化和可复现性，成为该领域最重要的基准之一。然而，OGB的节点属性主要为低维文本向量，丢失了论文原始文本中的丰富语义信息，更无法体现论文中的视觉信息（如图表、公式等），这限制了多模态图学习方法的研究和发展。

随着多模态学习技术的快速发展\cite{baltruvsaitis2019multimodal}，研究者们逐渐意识到，现实世界中的图数据往往同时包含多种模态的节点属性。多模态属性图（Multimodal Attributed Graph, MAG）正是这样一种数据结构，其节点不仅具有传统的标量或向量属性，还包含文本、图像等多种模态的信息。Yan等人\cite{yan2025magb}首次系统地定义了MAG学习范式，并构建了MAGB基准，包含覆盖电商（Amazon Movies、Amazon Toys、Amazon Grocery）和社交（Reddit）领域的5个数据集，每个节点都具有文本和图像双模态属性。MAGB的实验揭示了重要发现：在社交网络中视觉特征往往比文本特征更重要，而在电商领域中文本特征的贡献更为显著；简单拼接多模态特征存在模态不平衡问题。Wan等人\cite{wan2026openmag}进一步扩展了MAG基准的规模，提出了OpenMAG，包含19个数据集，覆盖6个领域，设计了8个评估任务，使用了16个编码器和24个基线模型。Zhu等人\cite{zhu2025mmgraph}提出了MM-GRAPH基准，包含5个多模态图数据集，系统地支持模态消融实验，能够量化各模态的独立和联合贡献。然而，这些MAG基准数据集均存在两个共性局限：一是仅包含文本和图像双模态，未涵盖公式这一学术论文特有的模态；二是领域覆盖局限于电商和社交领域，未涉及学术论文这一天然的多模态属性图场景。

学术论文作为科学知识的主要载体，天然构成了一个丰富的多模态属性图。在这一图中，每个节点代表一篇论文，节点间的有向边表示引用关系，而每个节点的属性则包含全文文本、实验图表、数学公式和表格等多种模态的信息。与电商和社交领域的多模态属性图相比，学术论文图具有其独特性：第一，公式是学术论文最核心的多模态成分之一，特别是在数学、物理和理论计算机科学领域，公式承载了大量结构化的数学信息，这些信息往往无法通过纯文本充分表达；第二，学术论文的引用网络具有明确的学术语义，引用关系反映了知识传播和学科发展的脉络；第三，学术图表（如架构图、实验结果图、流程图）与电商商品图或社交媒体图片在内容上有本质差异，需要领域适配的视觉编码方法。

现有工作中，面向学术论文领域的图基准数据集仍存在显著不足。在学术图基准方面，OGB\cite{hu2020ogb}和CS-TAG\cite{yan2023cstag}虽然覆盖了学术论文领域并提供了引用网络结构，但节点属性仅限于文本，缺乏图像和公式等多模态信息。经典的小规模学术图数据集（如Cora、CiteSeer、PubMed\cite{sen2008collective}）更是使用词袋（Bag-of-Words）表示作为节点特征，模态信息更为匮乏。在多模态学术文档基准方面，DocGenome\cite{xia2024docgenome}虽然提供了涵盖文本、图像、公式和布局四种模态的大规模科学文档数据集（50万篇论文），但仅关注单文档级别的多模态理解，未构建文档间的引用网络图结构。M3SciQA\cite{li2024m3sciqa}和MCiteBench\cite{hu2025mcitebench}分别关注多文档多模态科学问答和多模态引用生成任务，同样未将论文组织为图结构进行图学习任务的评估。

基于上述分析，目前学术界尚缺乏一个同时满足以下条件的数据集和基准：（1）面向学术论文领域；（2）包含文本、图像、公式等多种模态的节点属性；（3）具备引用网络图结构；（4）提供自动化数据采集与构建系统。为了填补这一空白，本文提出了AMMAG（Academic Multimodal Attributed Graph）——首个面向学术论文领域的多模态属性图Benchmark数据集及其自动化构建系统。

\section{国内外研究现状}

本文的研究涉及图神经网络与图基准、多模态属性图基准、多模态学术文档理解以及学术文献图构建与自动化采集四个交叉领域。近年来，GCN\cite{kipf2017gcn}、GAT\cite{velickovic2018gat}、GraphSAGE\cite{hamilton2017graphsage}和GIN\cite{xu2019gin}等模型推动了图学习的快速发展，OGB\cite{hu2020ogb}等基准提供了标准化评估框架，但现有基准的节点属性主要为低维文本向量，缺乏多模态信息。MAGB\cite{yan2025magb}、OpenMAG\cite{wan2026openmag}和MM-GRAPH\cite{zhu2025mmgraph}开创了多模态属性图学习范式，但均仅含文本和图像双模态且领域局限于电商和社交。DocGenome\cite{xia2024docgenome}等大规模科学文档基准证明了学术论文包含丰富的多模态信息，但未构建文档间的引用网络图结构。在自动化采集方面，AutoData\cite{ma2025autodata}等多智能体系统框架为复杂数据采集提供了技术基础\cite{guo2024llm_agent_survey}，但未针对学术MAG的自动化构建设计端到端流水线。第二章将对上述四个方面进行详细综述。

\subsection{现有工作总结}

% --- 用一张核心对比表总结 ---

\begin{table}[htbp]
  \centering
  \caption{现有工作特性对比}
  \label{tab:related_work_comparison}
  \begin{tabular}{lccccccc}
    \toprule
    特性 & OGB & CS-TAG & DocGenome & MAGB & OpenMAG & MM-GRAPH & \textbf{AMMAG} \\
    \midrule
    学术论文领域 & \checkmark & \checkmark & \checkmark & & & & \textbf{\checkmark} \\
    节点多模态属性 & & & \checkmark & \checkmark & \checkmark & \checkmark & \textbf{\checkmark(3种)} \\
    图结构(引用网络) & \checkmark & \checkmark & & \checkmark & \checkmark & \checkmark & \textbf{\checkmark} \\
    公式模态 & & & \checkmark & & & & \textbf{\checkmark} \\
    模态消融实验 & & & & \checkmark & \checkmark & \checkmark & \textbf{\checkmark(含公式)} \\
    自动化采集系统 & & & & & & & \textbf{\checkmark} \\
    \bottomrule
  \end{tabular}
\end{table}

如表\ref{tab:related_work_comparison}所示，目前没有任何一个工作同时满足学术论文领域、多模态节点属性（含公式模态）、引用网络图结构和自动化构建系统这四个条件。OGB和CS-TAG虽然覆盖了学术领域并提供了图结构，但缺乏多模态属性；DocGenome虽然提供了丰富的多模态信息，但缺乏图结构；MAGB、OpenMAG和MM-GRAPH虽然同时具备多模态属性和图结构，但其领域局限于电商和社交，且不包含公式模态。AMMAG正是填补这一空白的首次尝试。更详细的对比分析见第二章表\ref{tab:related_work_final}。

\section{本文贡献}

针对上述研究空白，本文提出了AMMAG——首个面向学术论文领域的多模态属性图Benchmark数据集及其基于多智能体的自动化构建系统。本文的主要贡献如下：

\textbf{（1）构建了首个学术多模态属性图Benchmark数据集AMMAG。}
AMMAG覆盖计算机科学、数学和物理学三个学科大类，包含cs.SI（社交信息学）、cs.LO（计算逻辑）、math.LO（数理逻辑）和physics.soc-ph（社会物理）四个子领域，共计45,672篇论文。每篇论文节点包含全文文本（Markdown格式）、实验图表（PNG格式）和数学公式（LaTeX格式）三种模态属性，并通过59,364条引用边构成引用网络图结构。与MAGB\cite{yan2025magb}和OpenMAG\cite{wan2026openmag}相比，AMMAG首次将公式模态引入MAG学习范式；与OGB\cite{hu2020ogb}和CS-TAG\cite{yan2023cstag}相比，AMMAG提供了原始多模态属性而非低维压缩向量。

\textbf{（2）设计了基于多智能体的自动化构建系统。}
本文设计并实现了包含五个专用智能体的自动化构建系统，采用编排器-智能体-工具三层架构。五个智能体分别负责论文下载、多模态提取、数据清洗、引用图构建和标准化格式化，实现了从arXiv元数据到标准化MAG数据集的端到端自动化流程。其中，多模态提取引擎基于DocLayout-YOLO\cite{zhao2024doclayoutyolo}实现了文本、图像和公式的并行提取；引用图构建模块通过双源引用提取和三级匹配引擎（精确匹配、子串匹配和模糊匹配）实现了自动化引用网络构建。与GROBID\cite{lopez2009grobid}等传统半自动工具相比，本系统支持多模态信息提取和完整的引用图自动构建。

\textbf{（3）进行了全面的Benchmark实验评估。}
本文在多模态消融、模型有效性、鲁棒性、Few-shot学习、链接预测和跨域迁移六组任务上，系统评估了13种模型（含7种标准GNN/MLP和6种多模态GNN）在7种模态组合下的性能表现，为学术MAG学习研究提供了全面的基线结果。

\textbf{（4）提供了首个学术多模态属性图的全面基线与关键发现。}
本文在六组任务上系统评估了13种模型与7种模态组合，产生了超过1,000组实验结果。关键发现包括：（a）T+F（文本+公式）是最佳模态组合（均值90.84\%），甚至略优于三模态T+I+F，公式比图像提供了更多互补增益——这是MAGB\cite{yan2025magb}等现有基准因缺少公式模态无法得出的独有发现；（b）不同学术领域对模态的偏好差异显著——cs.SI从多模态融合获益最大（+2.11pp），math.LO几乎不受图像影响，这种差异与领域模态分布直接相关；（c）公式模态的密度效应揭示了通用编码器的局限，表明需要专门的公式编码器；（d）跨域迁移增益接近零（$\pm$0.5pp），领域间子分类模式差异大，领域特定训练是必要的。这些发现为未来学术MAG学习中的多模态融合策略和编码器设计提供了重要的实证依据。

\section{论文组织结构}

本文共分为六章，各章内容安排如下：

第一章为绪论。介绍了本文的研究背景与意义，从图神经网络与Benchmark的发展、多模态属性图的兴起和学术文献的特殊价值三个角度阐述了研究动机。综述了国内外研究现状，分析了现有工作的局限性，最后阐述了本文的四点贡献。

第二章为相关工作。分别从图神经网络与图基准、多模态属性图基准、多模态学术文档理解、学术文献图构建与自动化采集四个方面详细综述了相关工作。通过对OGB、MAGB、DocGenome等代表性工作的深入分析，明确了AMMAG的研究定位。

第三章为系统架构设计。详细阐述了基于多智能体的自动化构建系统的三层架构设计，包括五个智能体的功能设计和关键技术实现。重点介绍了多模态提取引擎和引用图构建模块的设计原理与实现细节。

第四章为数据集构建与分析。介绍了AMMAG数据集的采集策略、统计分析和数据质量评估。通过与现有数据集的系统对比，展示了AMMAG的独特价值。

第五章为Benchmark实验与分析。在多模态消融、模型有效性、鲁棒性、Few-shot学习、链接预测和跨域迁移六组任务上系统评估了基线模型的性能，重点分析了模态消融实验揭示的学术领域特性。

第六章为总结与展望。总结了本文的主要工作和贡献，讨论了存在的不足，并展望了未来的研究方向。

% 第二章 相关工作
%%
% 第二章 相关工作（~7页）
%%

\chapter{相关工作}

本章从图神经网络与图基准、多模态属性图基准、多模态学术文档理解、学术文献图构建与自动化采集四个方面，系统综述与本文研究密切相关的工作，分析现有工作的局限性，明确本文的研究定位。

\section{图神经网络与图基准}

\subsection{图神经网络基础}

图神经网络是一类专门处理图结构数据的深度学习方法。给定一个图$G=(V, E)$，其中$V$为节点集合，$E$为边集合，每个节点$v_i \in V$具有特征向量$\mathbf{x}_i$。GNN的核心思想是通过消息传递（Message Passing）机制聚合邻居节点的信息来更新节点表示\cite{gilmer2017mpnn}。具体而言，在第$k$层迭代中，节点$v_i$的隐藏表示$\mathbf{h}_i^{(k)}$通过以下方式更新：
\begin{equation}
  \mathbf{h}_i^{(k)} = \text{UPDATE}^{(k)}\left(\mathbf{h}_i^{(k-1)}, \text{AGGREGATE}^{(k)}\left(\{\mathbf{h}_j^{(k-1)} : j \in \mathcal{N}(i)\}\right)\right)
  \label{eq:message_passing}
\end{equation}
其中$\mathcal{N}(i)$表示节点$v_i$的邻居集合，AGGREGATE函数负责聚合邻居信息，UPDATE函数负责结合自身信息和聚合后的邻居信息来更新节点表示。不同的AGGREGATE和UPDATE函数设计对应了不同的GNN模型。

\textbf{图卷积网络（GCN）。}Kipf和Welling\cite{kipf2017gcn}提出了图卷积网络，通过对图拉普拉斯矩阵的一阶切比雪夫多项式近似实现了高效的谱域图卷积。GCN的逐层传播规则为：
\begin{equation}
  \mathbf{H}^{(l+1)} = \sigma\left(\tilde{\mathbf{D}}^{-\frac{1}{2}}\tilde{\mathbf{A}}\tilde{\mathbf{D}}^{-\frac{1}{2}}\mathbf{H}^{(l)}\mathbf{W}^{(l)}\right)
  \label{eq:gcn}
\end{equation}
其中$\tilde{\mathbf{A}} = \mathbf{A} + \mathbf{I}$为添加自连接的邻接矩阵，$\tilde{\mathbf{D}}$为其度矩阵，$\mathbf{W}^{(l)}$为可学习参数，$\sigma$为非线性激活函数。GCN通过归一化的邻接矩阵实现了对称归一化的邻居聚合，其计算复杂度与图中边数呈线性关系，在大规模图上具有良好的可扩展性。GCN在引文网络分类等半监督节点分类任务上取得了当时的最佳性能，成为GNN领域最具影响力的工作之一。

\textbf{图注意力网络（GAT）。}Veli\v{c}kovi\'{c}等人\cite{velickovic2018gat}指出，GCN对所有邻居节点使用固定的聚合权重，无法自适应地区分不同邻居的重要性。为此，GAT引入了注意力机制，为每个邻居节点动态学习一个权重系数。注意力系数的计算方式为：
\begin{equation}
  \alpha_{ij} = \frac{\exp\left(\text{LeakyReLU}\left(\mathbf{a}^T[\mathbf{W}\mathbf{h}_i \| \mathbf{W}\mathbf{h}_j]\right)\right)}{\sum_{k \in \mathcal{N}(i)} \exp\left(\text{LeakyReLU}\left(\mathbf{a}^T[\mathbf{W}\mathbf{h}_i \| \mathbf{W}\mathbf{h}_k]\right)\right)}
  \label{eq:gat}
\end{equation}
其中$\mathbf{a}$为可学习的注意力向量，$\|$表示向量拼接。GAT还支持多头注意力机制，通过多组独立的注意力计算捕捉不同子空间的信息。注意力机制使GAT能够自适应地为重要邻居分配更高的权重，在存在噪声边的图上表现出更强的鲁棒性。

\textbf{GraphSAGE。}Hamilton等人\cite{hamilton2017graphsage}提出了GraphSAGE，解决了一个重要的实际问题——归纳式学习。与需要整图参与训练的GCN和GAT不同，GraphSAGE通过对邻居进行采样和聚合，能够对训练时未见过的节点生成表示。GraphSAGE支持多种聚合器（均值聚合、LSTM聚合、最大池化聚合），为每个节点采样固定数量的邻居，避免了计算开销随图规模增长的问题。

\textbf{图同构网络（GIN）。}Xu等人\cite{xu2019gin}从理论上分析了GNN的表达能力上界。他们证明，基于邻居聚合的GNN模型的表达能力不会超过Weisfeiler-Lehman（WL）图同构测试。基于这一理论分析，GIN采用求和聚合策略（而非均值或最大值聚合），使其在理论上达到最大的表达能力：
\begin{equation}
  \mathbf{h}_i^{(k)} = \text{MLP}^{(k)}\left((1 + \epsilon^{(k)}) \cdot \mathbf{h}_i^{(k-1)} + \sum_{j \in \mathcal{N}(i)} \mathbf{h}_j^{(k-1)}\right)
  \label{eq:gin}
\end{equation}
其中$\epsilon^{(k)}$为可学习的标量参数。

除上述经典模型外，Defferrard等人\cite{defferrard2016chebnet}提出的ChebNet通过切比雪夫多项式近似实现了快速的局部谱域滤波；Monti等人\cite{monti2017monet}提出的MoNet通过高斯混合模型定义了统一的空域聚合框架。Gilmer等人\cite{gilmer2017mpnn}提出了消息传递神经网络（MPNN）框架，将上述模型统一为消息传递的范式，为后续研究提供了清晰的理论框架。这些模型的不断发展推动了GNN在节点分类、图分类、链接预测等任务上的性能持续提升\cite{zhou2020gnn_survey, wu2021gnn_survey}。

\subsection{图基准数据集的发展}

高质量的基准数据集是推动GNN研究发展的关键基础设施。在图机器学习领域，基准数据集的发展经历了从经典小规模数据集到大规模标准化基准的演进过程。

\textbf{经典引文网络数据集。}Sen等人\cite{sen2008collective}系统整理了Cora、CiteSeer和PubMed三个引文网络数据集，长期以来被作为GNN模型评估的标准基准。Cora数据集包含2,708篇机器学习论文和5,429条引用边，节点属性为1,433维的词袋向量，标签为7个研究领域。CiteSeer数据集包含3,327篇论文和4,732条引用边，节点属性为3,703维词袋向量。PubMed数据集规模稍大，包含19,717篇糖尿病相关论文和44,338条引用边。这三个数据集在GNN研究的早期阶段发挥了重要作用，GCN\cite{kipf2017gcn}、GAT\cite{velickovic2018gat}等经典模型均在这些数据集上验证了其有效性。然而，这些数据集规模偏小，节点属性仅为低维词袋表示，难以充分评估模型在大规模图和丰富属性场景下的性能。

\textbf{开放图基准（OGB）。}Hu等人\cite{hu2020ogb}提出的OGB是图机器学习领域最重要的基准之一。OGB提供了涵盖节点属性预测、链接属性预测和图属性预测三大类任务的标准化评估框架。在节点分类任务中，ogbn-arxiv数据集包含约169,343篇arXiv论文节点和1,166,243条引用边，覆盖40个学术类别。ogbn-arxiv将论文的标题和摘要通过skip-gram模型编码为128维向量作为节点特征。ogbn-mag数据集则是一个异构学术图，包含约1.9M个节点（论文、作者、机构、研究领域四种类型）和约21M条边。OGB的重要贡献在于提供了统一的评估协议、自动化的性能排行榜和标准化的数据加载接口，极大地促进了图学习研究的可复现性和可比性。

然而，OGB在节点属性的丰富度上存在明显不足。ogbn-arxiv的128维skip-gram向量丢失了论文原始文本中的大量语义信息（如上下文语义、专业术语等），更无法体现论文中的视觉信息（实验图表、数学公式等）。这种信息损失限制了多模态图学习方法的研究和发展。

\textbf{文本属性图基准。}为了弥补OGB在文本信息方面的不足，Yan等人\cite{yan2023cstag}提出了CS-TAG基准，将论文的原始标题和摘要作为文本属性图（Text-Attributed Graph, TAG）的节点特征。CS-TAG包含Coauthor-CS和Coauthor-Phy两个数据集，支持在原始文本上评估不同的文本编码策略对图学习性能的影响。CS-TAG的实验表明，使用预训练语言模型（如SciBERT\cite{beltagy2019scibert}）编码文本可以显著提升节点分类性能。然而，CS-TAG仍局限于文本单模态，未引入图像和公式等其他模态信息。

Li等人\cite{li2024tegdb}提出了TEG-DB基准，首次将文本属性扩展到图的边上，支持边文本图学习任务。TEG-DB包含多个数据集，涵盖了论文引用（边文本为引用上下文）、社交互动（边文本为对话内容）等多种场景。然而，TEG-DB同样未引入多模态节点属性。

综上所述，现有图基准数据集在节点属性的丰富度上仍存在不足，特别是缺乏多模态节点属性（如图像、公式等）。表\ref{tab:graph_benchmarks}对主要图基准数据集的特性进行了系统对比。

\begin{table}[htbp]
  \centering
  \caption{主要图基准数据集特性对比}
  \label{tab:graph_benchmarks}
  \begin{tabular}{lccccc}
    \toprule
    数据集 & 节点数 & 节点属性 & 模态类型 & 图结构 & 规模 \\
    \midrule
    Cora\cite{sen2008collective} & 2,708 & 词袋(1433d) & 文本 & 引用网络 & 小 \\
    CiteSeer\cite{sen2008collective} & 3,327 & 词袋(3703d) & 文本 & 引用网络 & 小 \\
    ogbn-arxiv\cite{hu2020ogb} & 169K & skip-gram(128d) & 文本向量 & 引用网络 & 大 \\
    CS-TAG\cite{yan2023cstag} & 18K & 原始文本 & 文本 & 合作网络 & 中 \\
    TEG-DB\cite{li2024tegdb} & 多 & 节点+边文本 & 文本 & 多种类型 & 中 \\
    \bottomrule
  \end{tabular}
\end{table}

\section{多模态属性图基准}

多模态属性图（Multimodal Attributed Graph, MAG）是图结构与多模态节点属性的结合。与仅使用单一模态节点特征的传统图学习不同，MAG学习旨在利用多模态信息提升图学习任务的性能。Baltru\v{s}aitis等人\cite{baltruvsaitis2019multimodal}系统综述了多模态学习的发展历程和技术挑战，指出多模态融合、对齐和表示学习是多模态学习的三个核心问题。在图学习领域，如何有效地融合多模态节点属性与图结构信息，是一个新兴且具有挑战性的研究方向。

\subsection{MAGB：首个多模态属性图基准}

Yan等人\cite{yan2025magb}提出的MAGB是多模态属性图学习领域的开创性工作。MAGB首次系统地定义了MAG学习范式，并构建了包含5个数据集的基准，覆盖电商领域（Amazon Movies、Amazon Toys、Amazon Grocery）和社交领域（Reddit）。每个数据集中的节点都具有文本描述和产品/帖子图像两种模态属性。在文本编码方面，MAGB使用Sentence-BERT\cite{reimers2019sbert}将文本编码为384维向量；在图像编码方面，使用CLIP\cite{radford2021clip}将图像编码为512维向量。多模态融合采用简单拼接策略，将文本和图像向量拼接为896维的联合表示。

MAGB系统评估了16个基线模型在节点分类和链接预测任务上的表现，揭示了若干重要发现：（1）不同领域对模态的偏好存在显著差异——社交网络（Reddit）中视觉特征的贡献大于文本特征，而电商领域（Amazon）则更偏好文本特征；（2）简单拼接多模态特征存在模态不平衡问题，图像模态的高维度可能在某些场景下引入噪声；（3）图结构信息在多模态设置下仍能提供额外的增益，但这种增益因数据集而异。

然而，MAGB存在三个明显的局限：第一，领域覆盖局限于电商和社交领域，未涉及学术论文这一天然的多模态属性图场景；第二，模态类型仅限于文本和图像，不包含公式这一学术论文特有的模态；第三，部分数据集（如Amazon系列）不包含天然的图结构，其图结构是基于用户共同购买或共同评论关系构建的，与引用网络等学术图结构有本质差异。

\subsection{OpenMAG：大规模综合基准}

Wan等人\cite{wan2026openmag}提出了OpenMAG，在规模和评估维度上大幅扩展了MAG基准。OpenMAG包含19个数据集，覆盖电商、社交、电影、音乐、学术引用和知识图谱等6个领域，设计了8个评估任务（包括节点分类、链接预测、图分类等），评估了16个编码器和24个基线模型。OpenMAG的一个重要贡献是提供了统一的评估框架，支持不同编码器组合和融合策略的系统比较。

OpenMAG的实验进一步验证了MAGB的发现：多模态信息在不同领域的贡献存在差异，且简单的拼接融合策略并非最优选择。然而，OpenMAG在模态类型上仍局限于文本和图像双模态，同样不包含公式模态。此外，虽然OpenMAG声称覆盖了学术引用领域，但其学术数据集（如ogbn-arxiv）实际上并不包含原始的多模态属性，仅使用了OGB提供的低维文本向量。

\subsection{MM-GRAPH：模态消融基准}

Zhu等人\cite{zhu2025mmgraph}提出了MM-GRAPH，专注于模态消融实验的系统评估。MM-GRAPH包含5个多模态图数据集，支持7种模态组合的消融实验：仅图结构（G）、仅文本（T）、仅视觉（V）、图+文本（GT）、图+视觉（GV）、文本+视觉（TV）和全部模态（GTV）。这种设计使得研究者能够量化各模态的独立贡献和交互效应。

MM-GRAPH的消融实验揭示了一个重要现象：在大多数数据集上，文本+视觉的组合优于任何单一模态，但加入图结构信息后的三模态组合（GTV）并不总是优于双模态组合（TV）。这一发现表明，模态之间的交互效应并非简单的叠加关系，需要设计更精细的融合策略。然而，MM-GRAPH的领域同样局限于电商和社交领域，模态类型仅为文本和视觉两种，未涵盖公式等其他模态。

\subsection{现有MAG基准的共性问题}

综合分析MAGB\cite{yan2025magb}、OpenMAG\cite{wan2026openmag}和MM-GRAPH\cite{zhu2025mmgraph}三大MAG基准，可以发现它们存在以下共性局限：

\textbf{（1）模态类型单一。}现有MAG基准均只包含文本和图像两种模态，未涵盖公式这一学术论文特有的模态。公式承载了大量的结构化数学信息，在学术文献的分类和理解中具有独特价值。缺少公式模态使得现有MAG基准无法评估公式在学术图学习任务中的贡献。

\textbf{（2）领域覆盖有限。}现有MAG基准的数据集主要来源于电商和社交领域。学术论文作为一个天然的多模态属性图场景，其节点（论文）具有全文文本、实验图表、数学公式和表格等多种模态属性，其边（引用关系）具有明确的学术语义。电商和社交领域的发现（如"社交偏好视觉、电商偏好文本"）是否适用于学术领域，目前尚无系统的研究。

\textbf{（3）图结构差异。}电商和社交领域的图结构多为用户-商品交互图或社交关系图，其边的密度和分布特征与学术论文的引用网络有本质差异。引用网络通常具有有向性、稀疏性和幂律度分布等特征，这些特征对GNN消息传递的效果有重要影响。

本文提出的AMMAG正是为了填补上述空白：首次将公式模态引入MAG基准，将领域覆盖扩展到学术论文领域，并基于真实的引用网络结构构建图数据。

\section{多模态学术文档理解}

学术文档的多模态理解是自然语言处理和文档智能的重要研究方向。学术论文作为一种特殊的多模态文档，同时包含文本、图表、公式和表格等多种模态信息，其理解需要综合运用计算机视觉、自然语言处理和信息提取等多种技术。

\subsection{DocGenome：大规模科学文档基准}

Xia等人\cite{xia2024docgenome}提出的DocGenome是目前规模最大的科学文档多模态基准。DocGenome包含约50万篇科学文档，涵盖文本、图像、数学公式和页面布局四种模态，并提供了13种细粒度的布局属性标注和完整的LaTeX源代码。DocGenome覆盖了153个学科领域，从计算机科学到物理学、数学、生物学等，为训练和评估多模态大语言模型提供了丰富的数据基础。

DocGenome在文档布局分析方面采用了层次化的标注体系，将页面元素分为标题、正文、图片、图片标题、表格、表格标题、行内公式、行间公式、脚注、页眉、页脚等13种类型。这种细粒度的标注使得模型能够学习不同页面元素之间的语义关系。DocGenome的构建流程包括PDF解析、布局检测、内容提取和质量控制四个阶段，其中布局检测使用了DocLayout-YOLO等先进的文档布局分析模型。

DocGenome与AMMAG的关系体现在以下方面：一方面，DocGenome证明了学术论文确实包含丰富的多模态信息（特别是公式模态），这与本文的研究动机一致；另一方面，DocGenome专注于单文档级别的多模态理解任务（如文档分类、信息提取、视觉问答等），未构建文档间的引用网络图结构。AMMAG则将多模态属性与图结构相结合，首次在引用网络图上评估多模态图学习方法的性能。

\subsection{多模态科学问答与引用生成}

Li等人\cite{li2024m3sciqa}提出了M3SciQA，一个多文档多模态科学问答基准。M3SciQA要求模型同时理解多篇科学论文中的文本、图表和公式信息来回答复杂问题，强调跨文档的推理能力。M3SciQA包含多种问题类型，包括数值推理、图表解读、跨文档比较等，能够全面评估模型在科学文档理解方面的能力。M3SciQA的实验表明，当前最先进的多模态大语言模型在涉及公式理解和跨文档推理的问题上表现仍然不佳，说明学术文档的多模态理解仍是一个具有挑战性的任务。

Hu等人\cite{hu2025mcitebench}提出了MCiteBench，一个多模态引用生成基准。MCiteBench要求模型根据论文的文本和图表信息生成带有引用标记的文本，评估模型在学术写作辅助方面的能力。MCiteBench关注的是引用文本的生成质量，而非引用网络的结构特征。

Roberts等人\cite{roberts2024scifibench}提出了SciFIBench，专注于科学图表理解任务的评估。SciFIBench包含多种类型的科学图表（折线图、柱状图、散点图、示意图等），评估模型解读图表内容、提取数据趋势和进行科学推理的能力。SciFIBench的实验表明，即使是最先进的多模态模型在解读复杂科学图表时仍有较大提升空间。

\subsection{文档理解的技术基础}

多模态学术文档理解的技术基础包括文档布局分析、光学字符识别（OCR）和公式识别等。在文档布局分析方面，DocLayout-YOLO\cite{zhao2024doclayoutyolo}通过多样化的合成数据和全局-局部自适应感知机制，在科学文档的版面检测上取得了优异的性能。LayoutLMv3\cite{huang2022layoutlmv3}和DiT\cite{li2022dit}则通过预训练策略学习了文档图像中文本、布局和视觉信息的联合表示。Blecher等人\cite{blecher2024nougat}提出的Nougat模型能够将学术PDF文档转换为结构化的Markdown文本，保留了公式、表格和图片的语义信息。

在公式识别方面，UniMERNet\cite{wang2024unimernet}提出了一个通用的数学公式识别网络，能够将文档中的数学公式图像转换为LaTeX代码，在真实场景的公式识别任务上表现优异。在OCR方面，PaddleOCR\cite{paddleocr}是一个开源的多语言OCR工具包，支持多种文字识别场景。本文的自动化构建系统正是基于DocLayout-YOLO、UniMERNet和PaddleOCR等技术实现了学术论文的多模态信息提取。

\subsection{本节小结}

本节综述的多模态学术文档理解工作在文档级别的多模态信息提取和理解方面取得了重要进展。然而，这些工作关注的是单文档或少量文档的理解任务（如问答、引用生成、图表解读），均未将大量论文组织为图结构进行图学习任务的评估。学术论文间的引用关系构成了一种天然的图结构，利用这一图结构进行消息传递可能有助于更好地理解论文的语义。AMMAG正是将学术论文的多模态属性嵌入引用网络图结构中，首次在学术多模态属性图上开展系统评估。

\section{学术文献图构建与自动化采集}

学术文献图的构建涉及PDF文档解析、元数据提取、引用网络构建等多个环节，是实现学术多模态属性图数据集的基础。本节综述相关工具、数据集和自动化采集系统。

\subsection{学术PDF解析工具}

学术PDF文档的解析是构建学术文献图的第一步。Lopez\cite{lopez2009grobid}开发的GROBID（GeneRation Of BIbliographic Data）是目前使用最广泛的学术PDF结构化解析工具。GROBID基于条件随机场（CRF）和深度学习模型，能够从PDF文档中提取标题、作者、摘要、参考文献等结构化信息。GROBID采用模块化的架构设计，包含头部提取、参考文献解析、全文提取等多个模块，可以通过REST API进行调用。然而，GROBID是一种半自动工具，不支持图像和公式等多模态信息的提取，也无法自动构建完整的引用网络图。

Tkaczyk等人\cite{tkaczyk2015cermine}提出了CERMINE，一个基于规则和机器学习的学术PDF内容提取系统。CERMINE能够自动提取文档的元数据（标题、作者、摘要）、正文内容和参考文献，其准确率在多个标准数据集上与GROBID相当。ScienceParse是Allen人工智能研究所（AI2）开发的PDF解析工具，基于深度学习模型提取论文的结构化信息。与GROBID相比，ScienceParse更加轻量化，但在复杂文档格式上的处理能力稍弱。

这些传统PDF解析工具的共同局限在于：只能提取文本类型的元数据和内容，不支持图像、公式等多模态信息的提取；引用关系的提取依赖参考文献列表的解析，无法自动与论文数据库进行匹配构建完整的引用网络图；多为半自动工具，需要人工干预和配置，难以实现大规模的端到端自动化。

\subsection{大规模学术数据集}

大规模学术数据集为学术文献图的研究提供了重要的数据基础。Lo等人\cite{lo2020s2orc}提出的S2ORC（Semantic Scholar Open Research Corpus）是目前最大的开放学术语料之一，包含约4100万篇论文、约13亿条引用关系和约380万篇全文文本。S2ORC为每篇论文提供了结构化的元数据（标题、作者、摘要、发表年份、引用列表等），并将论文间的引用关系组织为引用网络图。然而，S2ORC主要提供文本类型的元数据，未提供论文的原始多模态内容（全文、图像、公式等）。

Tang等人\cite{tang2008arnetminer}开发的AMiner是一个学术搜索引擎和挖掘平台，提供了超过2亿篇论文的元数据搜索服务。AMiner构建了作者-论文-会议-机构的异构学术网络，支持学术关系挖掘和专家发现等任务。Priem等人\cite{priem2022openalex}提出的OpenAlex是一个开源的学术元数据仓库，提供了超过2.5亿篇学术作品的元数据和引用关系，覆盖论文、数据集、软件等多种学术成果类型。Clement等人\cite{clement2019arxiv}对arXiv数据集进行了系统分析，指出arXiv包含超过150万篇论文、约670万条引用边和110亿词的全文语料，是研究学术文献的重要数据源。

然而，上述数据集主要提供论文的元数据信息（标题、摘要、作者、引用关系等），未提供论文的原始多模态内容。研究者在使用这些数据集进行图学习研究时，只能使用低维的文本向量或词袋表示作为节点特征，丢失了大量有价值的多模态信息。

\subsection{自动化数据采集与多智能体系统}

随着大语言模型（LLM）技术的发展，基于多智能体系统的自动化数据采集成为可能。Ma等人\cite{ma2025autodata}提出的AutoData是一个基于多智能体系统的开放网络数据采集框架。AutoData利用多个大语言模型智能体协同完成数据的搜索、提取、验证和整合，实现了高度自动化的数据采集流程。AutoData的核心设计包括：搜索智能体负责从网络中发现相关数据源；提取智能体负责从网页中提取结构化数据；验证智能体负责检查数据的完整性和一致性；整合智能体负责将不同来源的数据合并为统一格式。然而，AutoData是一个通用框架，未针对学术多模态属性图的自动化构建进行专门设计。

Guo等人\cite{guo2024llm_agent_survey}系统综述了基于大语言模型的多智能体系统研究进展，将多智能体系统分为协作型、对抗型和层次型三种类型。在数据采集和处理的场景中，协作型多智能体系统最为适用——多个专用智能体各司其职，通过编排器协调执行顺序和数据流。本文的系统设计正是采用了这种协作型多智能体架构。

在图学习框架方面，Fey和Lenssen\cite{fey2019pyg}开发了PyTorch Geometric（PyG），提供了高效的GNN模型实现和数据处理工具。Wang等人\cite{wang2019dgl}开发了Deep Graph Library（DGL），支持多种深度学习框架后端。这些框架为图学习实验的实现提供了技术基础，本文的Benchmark实验即基于PyG框架实现。

\subsection{本节小结}

综合以上分析，现有学术文献图构建工具和数据采集系统存在以下不足：传统PDF解析工具（如GROBID）不支持多模态信息提取和引用网络自动构建；大规模学术数据集（如S2ORC）仅提供元数据，不包含原始多模态内容；通用多智能体框架（如AutoData）未针对学术MAG的自动化构建设计端到端流水线。本文提出的基于多智能体的自动化构建系统正是为了解决这些问题，实现了从arXiv元数据到标准化MAG数据集的端到端自动化流程。

\section{本章小结}

本章从四个方面系统综述了与本文研究密切相关的工作。在图神经网络与图基准方面，GCN\cite{kipf2017gcn}、GAT\cite{velickovic2018gat}、GraphSAGE\cite{hamilton2017graphsage}和GIN\cite{xu2019gin}等模型为图学习提供了坚实的方法基础，OGB\cite{hu2020ogb}等基准为评估提供了标准化框架，但现有基准在节点属性的丰富度上存在不足。在多模态属性图基准方面，MAGB\cite{yan2025magb}、OpenMAG\cite{wan2026openmag}和MM-GRAPH\cite{zhu2025mmgraph}开创了MAG学习范式，但均不包含公式模态且未覆盖学术论文领域。在多模态学术文档理解方面，DocGenome\cite{xia2024docgenome}等基准证明了学术论文包含丰富的多模态信息，但未将论文组织为图结构进行图学习评估。在自动化采集方面，GROBID\cite{lopez2009grobid}等工具和AutoData\cite{ma2025autodata}等系统提供了技术基础，但未针对学术MAG设计端到端流水线。

\begin{table}[htbp]
  \centering
  \caption{现有工作特性综合对比}
  \label{tab:related_work_final}
  \begin{tabular}{p{2.8cm}ccccccccc}
    \toprule
    特性 & \rotatebox{70}{OGB} & \rotatebox{70}{CS-TAG} & \rotatebox{70}{DocGenome} & \rotatebox{70}{MAGB} & \rotatebox{70}{OpenMAG} & \rotatebox{70}{MM-GRAPH} & \rotatebox{70}{S2ORC} & \rotatebox{70}{AutoData} & \rotatebox{70}{\textbf{AMMAG}} \\
    \midrule
    学术论文领域 & \checkmark & \checkmark & \checkmark & & & & \checkmark & & \textbf{\checkmark} \\
    节点多模态属性 & & & \checkmark & \checkmark & \checkmark & \checkmark & & & \textbf{\checkmark(3种)} \\
    图结构(引用网络) & \checkmark & \checkmark & & \checkmark & \checkmark & \checkmark & \checkmark & & \textbf{\checkmark} \\
    公式模态 & & & \checkmark & & & & & & \textbf{\checkmark} \\
    模态消融实验 & & & & \checkmark & \checkmark & \checkmark & & & \textbf{\checkmark(含公式)} \\
    自动化采集系统 & & & & & & & & \checkmark & \textbf{\checkmark} \\
    \bottomrule
  \end{tabular}
\end{table}

如表\ref{tab:related_work_final}所示，通过以上分析可以得出结论：目前没有任何一个工作同时满足学术论文领域、多模态节点属性（含公式模态）、引用网络图结构和自动化构建系统这四个条件。AMMAG正是填补这一空白的首次尝试，它继承了OGB和CS-TAG的学术领域和图结构优势，借鉴了DocGenome的多模态信息提取技术，融合了MAGB的MAG学习范式，并通过基于多智能体的自动化构建系统实现了端到端的数据集生成。

% 第三章 系统架构设计
%%
% 第三章 系统架构设计（~10页）← 重点章
%%

\chapter{系统架构设计}

本章详细阐述基于多智能体系统（Multi-Agent System, MAS）\cite{guo2024llm_agent_survey, wu2023autogen}的自动化构建系统设计。多智能体系统是近年来随着大语言模型技术快速发展而兴起的系统架构范式，其核心思想是将复杂任务分解为多个专用智能体协同完成。与传统的单体系统相比，多智能体系统具有职责分离、可扩展性强和容错性好等优势，特别适合数据构建这类涉及多个异构处理阶段（PDF下载、多模态提取、文本清洗、引用图构建等）的复杂流程。AMMAG的自动化构建系统正是采用了这一架构范式，通过五个专用智能体的协同工作，实现了从arXiv元数据到标准化多模态属性图Benchmark数据集的端到端自动化。本章首先介绍系统的总体架构和设计原则，然后依次描述五个专用智能体的功能设计、核心算法实现和性能优化策略。

\section{系统总体架构}

\subsection{设计目标与原则}

AMMAG自动化构建系统基于多智能体系统架构，其设计目标是实现从arXiv元数据到标准化多模态属性图数据集的端到端自动化流程。学术Benchmark数据集的构建涉及多个异构处理阶段——从大规模元数据过滤、PDF批量下载、多模态内容提取、文本规范化清洗，到引用关系的提取与匹配、最终的数据集格式化——这些阶段在数据格式、计算资源需求和处理策略上差异显著。传统的单体系统难以同时满足这些异构需求，而多智能体系统通过将每个处理阶段封装为一个专用智能体，使每个智能体可以针对其特定任务进行独立优化。具体而言，系统需要满足以下四项设计目标：（1）端到端自动化——从原始arXiv论文列表到可直接用于图学习实验的标准化数据集，无需人工干预；（2）多模态提取——同时提取论文PDF中的文本、图像和数学公式三种模态信息；（3）引用图构建——自动匹配论文间的引用关系，构建引用网络图结构；（4）可扩展性——支持新增学术领域、新模态类型和新处理工具的灵活扩展。

基于上述目标，系统遵循以下四项设计原则：

\textbf{智能体职责分离原则。}这是多智能体系统最核心的设计原则。将复杂的Benchmark数据构建流程分解为多个功能独立的智能体，每个智能体专注于单一数据处理阶段。例如，论文下载智能体只需关注PDF获取，多模态提取智能体只需关注内容提取，引用图构建智能体只需关注引用关系匹配。职责分离使得每个智能体可以独立开发、测试和优化，降低了系统的整体复杂度。当某个阶段需要升级（如替换更先进的公式识别模型）时，只需修改对应的智能体，不影响其他智能体的功能。

\textbf{可扩展原则。}通过工具注册机制（MCP工具管理器）将底层工具与上层智能体解耦。新增处理工具只需实现标准化的工具接口并注册到工具管理器中，无需修改智能体或编排器的代码。这种设计使得多智能体系统可以灵活地扩展新的处理能力——例如，未来新增表格结构提取功能只需注册一个新的工具接口。

\textbf{智能体间松耦合原则。}智能体之间通过标准化的数据接口（任务配置字典和上游数据字典）进行通信，而非直接调用。每个智能体接收任务配置和上游智能体的输出作为输入，产生自己的输出传递给下游。这种松耦合设计使得智能体可以被替换、重排或并行执行，而不影响其他智能体的功能。

\textbf{确定性原则。}数据清洗和引用匹配等关键步骤采用纯规则驱动的方法实现，不依赖大语言模型，确保相同的输入始终产生相同的输出，保证数据处理的可复现性。这一原则确保了多智能体系统在数据处理层面的可靠性。

\subsection{三层架构设计}

系统采用编排器-智能体-工具三层多智能体架构，如图\ref{fig:architecture}所示。这一架构借鉴了AutoGen\cite{wu2023autogen}等多智能体框架的设计思想，将复杂的Benchmark构建任务组织为"编排器协调——智能体执行——工具支撑"的层次化协同模式。

\begin{figure}[htbp]
  \centering
  \includegraphics[width=0.95\textwidth]{images/architecture.png}
  \caption{基于多智能体系统的三层总体架构图}
  \label{fig:architecture}
\end{figure}

\textbf{编排器层}（Orchestrator Layer）位于多智能体系统的最上层，充当中央协调器的角色，负责调度五个智能体的执行顺序、管理智能体间的数据流转和处理异常情况。在多智能体系统架构中，编排器的存在是必要的——五个智能体各自独立运行，需要一个全局协调者来确保它们按照正确的顺序执行并传递中间结果。系统提供两种编排策略：顺序编排器（\texttt{DataCrawlerOrchestrator}）按固定顺序依次调度各智能体，适用于小规模数据处理；并行编排器（\texttt{ParallelDataCrawlerOrchestrator}）支持多GPU并行提取和多进程并行清洗，通过配置参数控制GPU数量和CPU工作进程数，适用于大规模数据处理。

\textbf{智能体层}（Agent Layer）是多智能体系统的核心层，包含五个功能专一、接口统一的专用智能体。多智能体系统的核心优势在于每个智能体可以针对其特定任务进行深度优化：（1）\texttt{ArxivBatchDownloaderAgent}——从arXiv元数据中按领域过滤论文，优化了大规模JSONL文件的流式读取；（2）\texttt{PaperCollectorAgent}——下载PDF并调用多模态提取引擎，优化了多GPU并行批量处理；（3）\texttt{DataCleanerAgent}——执行六步清洗流水线，采用纯规则驱动确保确定性；（4）\texttt{ReferenceGraphBuilderAgent}——构建引用网络图，设计了倒排索引和三级匹配引擎；（5）\texttt{DatasetFormatterAgent}——格式化为标准化数据集。每个智能体继承自抽象基类\texttt{BaseAgent}，实现统一的\texttt{execute()}接口，接收任务配置和上游智能体的输出数据，产生自己的输出传递给下游智能体。这种统一的接口设计是多智能体系统实现灵活编排的基础——编排器无需了解每个智能体的内部实现，只需调用其\texttt{execute()}方法并传递数据即可。

\textbf{工具层}（Tool Layer）位于多智能体系统的最下层，由\texttt{MCPToolManager}统一管理，为智能体层提供标准化的工具支撑。工具管理器采用插件式注册模式，将底层工具（PDF下载、多模态提取、数据清洗、格式转换等）封装为标准化的工具接口。智能体通过\texttt{call\_mcp\_tool()}方法调用工具，实现了智能体与具体工具实现的解耦。这种设计使得多智能体系统中的智能体无需关心底层工具的实现细节，降低了智能体的开发复杂度。

\subsection{数据流全流程}

系统的数据流如图\ref{fig:pipeline}所示。在多智能体系统架构下，完整的数据处理流程被分解为五个阶段，每个阶段由一个专用智能体负责，智能体之间通过标准化的数据接口传递中间结果：

\begin{figure}[htbp]
  \centering
  \includegraphics[width=0.95\textwidth]{images/pipeline.png}
  \caption{数据构建Pipeline全流程图}
  \label{fig:pipeline}
\end{figure}

（1）\texttt{ArxivBatchDownloaderAgent}从arXiv元数据集（约299万篇论文的JSON Lines文件）中按目标领域（如cs.SI）过滤论文，按更新日期排序后输出论文列表。

（2）\texttt{PaperCollectorAgent}接收论文列表，批量下载PDF文件，然后调用多模态提取引擎从每篇PDF中提取文本（Markdown格式）、图像（PNG格式）和公式（LaTeX格式）。

（3）\texttt{DataCleanerAgent}对提取的Markdown文本执行六步清洗流水线，修复PDF转换引入的格式问题，提取和格式化参考文献。

（4）\texttt{ReferenceGraphBuilderAgent}从PDF和Markdown双源提取引用信息，通过三级匹配引擎将引用与arXiv元数据中的论文进行匹配，构建引用网络图。

（5）\texttt{DatasetFormatterAgent}将清洗后的多模态数据、引用边信息和论文元数据整合为标准化的多模态属性图数据集。

\section{ArxivBatchDownloaderAgent}

ArxivBatchDownloaderAgent负责从大规模arXiv元数据集中按学术领域筛选目标论文。该智能体的核心功能是从约299万篇论文的元数据中高效地过滤出指定领域（如cs.SI）的论文，并按更新日期排序后输出给下游智能体。

\subsection{元数据过滤算法}

arXiv元数据集以JSON Lines格式存储，每行包含一篇论文的元数据信息（标题、作者、领域标签、摘要、更新日期等）。由于文件总大小约5.3GB，无法一次性加载到内存中。因此，系统采用流式读取策略，逐行解析JSON记录并按条件过滤。

过滤算法的核心逻辑为：对每条元数据记录，提取其\texttt{categories}字段的第一个标签作为主领域标签（primary category），判断是否与目标领域匹配。例如，一篇论文的\texttt{categories}字段为\texttt{"cs.SI cs.LG"}，则其主领域标签为\texttt{cs.SI}。通过仅匹配主领域标签，系统确保了筛选结果的领域纯度，避免了跨领域论文的引入。

匹配成功的论文按\texttt{update\_date}字段降序排序，优先选择最新更新的论文，然后取前\texttt{limit}篇作为最终结果。以cs.SI领域为例，系统从元数据中筛选出12,000篇候选论文，最终成功下载11,951篇（下载成功率99.6\%）。

\section{PaperCollectorAgent与多模态提取引擎}

PaperCollectorAgent是系统中计算密集度最高的智能体，负责批量下载论文PDF并调用多模态提取引擎从中提取文本、图像和公式三种模态信息。

\subsection{多模态提取引擎}

多模态提取引擎基于PDF-Extract-Kit工具包\cite{zhao2024doclayoutyolo}实现，其处理流程如图\ref{fig:extraction}所示。对于每篇论文PDF，引擎首先使用DocLayout-YOLO\cite{zhao2024doclayoutyolo}进行版面检测，识别出页面中的五种区域类型：文本块、图片、行内公式、行间公式和表格。

\begin{figure}[htbp]
  \centering
  \includegraphics[width=0.95\textwidth]{images/extraction.png}
  \caption{多模态提取引擎处理流程}
  \label{fig:extraction}
\end{figure}

版面检测完成后，引擎根据区域类型分别进行处理：

\textbf{文本区域：}使用PaddleOCR\cite{du2020ppocr}对文本区域进行光学字符识别（OCR），将识别结果按段落结构组织为Markdown格式。PaddleOCR是一个开源的超轻量级OCR系统，支持多语言文字识别，在中文和英文文档上均有良好的识别精度。

\textbf{图片区域：}直接从PDF中裁剪图片区域的像素数据，保存为PNG格式。每篇论文的所有图片存储在独立的\texttt{images/}子目录中。

\textbf{公式区域：}首先使用YOLOv8检测行内和行间公式区域，然后使用UniMERNet\cite{wang2024unimernet}将公式图像识别为LaTeX代码。UniMERNet是一个通用的数学公式识别网络，能够处理复杂的手写和印刷体数学表达式，包括分数、根号、矩阵、积分等高级数学结构。

每篇论文的提取结果以统一格式存储：\texttt{paper.md}（Markdown全文）、\texttt{paper.npy}（结构化数据文件）和\texttt{images/}（图片目录）。

\subsection{多GPU并行提取}

多模态提取是整个系统中计算开销最大的环节。每次启动提取模型（DocLayout-YOLO、UniMERNet等）需要约30秒的模型加载时间。如果逐篇处理，仅模型加载就需消耗大量时间。

系统设计了\texttt{GPUExtractionPool}实现多GPU并行批量提取。其核心优化策略如下：

\textbf{批量处理。}将多篇论文的PDF通过符号链接组织到临时目录中，使用\texttt{-{}-pdf\_dir}参数一次性提交给提取引擎。这样，模型只需加载一次即可处理整个批次的论文（默认批次大小为50篇），将每篇论文的模型加载开销从30秒降至约0.6秒。

\textbf{多GPU并行。}系统将待处理的论文列表按轮转（round-robin）方式分配到多个GPU上，每个GPU由独立的线程驱动。默认配置使用8个GPU，通过\texttt{CUDA\_VISIBLE\_DEVICES}环境变量隔离每个GPU的可见设备，确保不同GPU上的提取进程互不干扰。

\textbf{失败回退。}批量提取中如果某篇PDF处理失败（如损坏的PDF文件），不会阻塞整个批次。系统对失败的论文回退到逐篇提取模式（fallback sequential），单独处理每篇论文，保证单篇论文的失败不影响批次中其他论文的处理。

\textbf{跳过已处理。}在开始提取前，系统检查每篇论文的输出目录是否已存在\texttt{.md}或\texttt{.npy}文件。如果已存在，则跳过该论文的提取，避免重复计算。

\section{DataCleanerAgent}

DataCleanerAgent负责对提取的Markdown文本进行规则驱动的清洗处理。该智能体不使用大语言模型，完全通过正则表达式和启发式规则实现，确保清洗过程的确定性和可复现性。

\subsection{清洗流水线}

清洗流水线由六个步骤组成，依次对Markdown文本进行规范化处理：

\textbf{步骤一：断词修复。}PDF转换过程中，英文单词经常被连字符和换行符分割。例如，"learning"可能被转换为"learn-\\ning"。断词修复步骤使用正则表达式\texttt{r"(\textbackslash w+)-\textbackslash s*\textbackslash n\textbackslash s*(\textbackslash w+)"}匹配被分割的单词模式，将断词部分重新合并为完整的单词。

\textbf{步骤二：段落合并。}PDF提取的文本经常出现段落被错误分行的情况。段落合并步骤逐行扫描文本，判断相邻两行是否应合并为同一段落。合并的条件包括：当前行不以句号、感叹号、问号等句末标点结尾；下一行不以Markdown标题（\texttt{\#}）、列表标记（\texttt{-}、\texttt{*}、\texttt{1.}）或引用编号（\texttt{[N]}）开头；下一行不以大写字母+逗号+首字母缩写模式开头（避免合并参考文献的作者列表）。

\textbf{步骤三：重复段落去除。}使用MD5哈希值对段落进行去重。对于长度超过50个字符的段落，计算其去除空白字符后的MD5哈希值，如果该哈希值已被记录则跳过该段落。长度不足50字符的短段落无条件保留，避免误删标题等短文本。

\textbf{步骤四：公式格式标准化。}统一LaTeX公式的格式规范：去除\texttt{\$\$}前后多余的空白字符；确保块级公式（\texttt{\$\$}）独占一行；统一行内公式（\texttt{\$...\$}）内部的前后空格。

\textbf{步骤五：参考文献提取与格式化。}这是最复杂的清洗步骤。系统使用三种策略依次尝试提取和格式化参考文献：（1）按\texttt{[N]}数字标签分割参考文献条目；（2）按作者姓名模式分割；（3）按年份边界分割。每种策略提取的参考文献条目都需通过有效性验证（包含4位年份、长度至少30字符、包含作者姓名模式或已知会议/期刊名称）。最终将参考文献按序编号为\texttt{[1], [2], ...}的标准格式。

\textbf{步骤六：冗余空白清理。}压缩三个及以上连续空行为两个空行，去除每行末尾的空白字符，确保文件以单个换行符结尾。

\subsection{并行清洗}

清洗步骤支持多进程并行执行。系统使用Python的\texttt{multiprocessing.Pool}创建指定数量的工作进程（默认8个），每个进程独立执行完整的六步清洗流水线。由于清洗过程是纯CPU操作且无共享状态，多进程并行可以实现近线性的加速比。

\section{ReferenceGraphBuilderAgent}

ReferenceGraphBuilderAgent是系统中技术复杂度最高的智能体，负责从论文中提取引用信息并构建引用网络图。该智能体采用双源提取策略获取原始引用信息，通过三级匹配引擎将引用与arXiv元数据中的论文进行匹配，最终输出引用边的CSV文件。

\subsection{双源引用提取}

学术论文的参考文献信息可以从两个数据源提取，两个源具有互补性：

\textbf{数据源一：PDF直接解析。}使用pdfplumber库直接解析PDF文件，提取参考文献页面的原始文本。该数据源的优势在于保留了PDF的精确排版信息，包括双栏布局。系统通过分析页面中字符的x坐标分布来检测双栏布局——在页面中间40\%的宽度范围内寻找最大的x坐标间隙，如果该间隙超过12磅（point），则判定为双栏布局，分别裁剪左右两栏进行文本提取。此外，系统通过检测跨页面重复的首行和末行来识别并移除页眉和页脚。

\textbf{数据源二：Markdown提取。}从已提取的Markdown文本中搜索参考文献章节。系统使用五种正则表达式模式匹配不同的章节标题格式（\texttt{\#\# References}、\texttt{\# References}、\texttt{References\\n}等）。对每个候选区域，系统计算一个综合评分函数，根据引用编号密度、arXiv标识符密度、年份密度等指标评估该区域作为参考文献章节的可能性。

两种数据源的融合显著提升了引用信息的完整性。在小规模验证实验中，对于cs.SI领域的论文样本，单源提取仅获得344条引用边，而双源融合后获得1,064条引用边，提升了209\%。

\subsection{多格式引用解析}

学术论文的参考文献采用多种不同的格式，系统支持识别和解析七种主要的引用格式类型：（1）方括号数字编号\texttt{[N]}格式；（2）圆括号数字编号\texttt{(N)}格式；（3）作者-年份编码\texttt{[Author+YY]}格式；（4）全角方括号数字编号格式；（5）点号编号\texttt{N.}格式；（6）空格编号\texttt{N Author}格式；（7）无编号格式（通过行级别或年份边界分割）。

PDF提取过程中，多条引用记录可能被错误地合并为一行（特别是当引用条目较短时）。系统设计了多信号边界检测算法来处理这种情况。该算法为四种边界信号设定权重：方括号编号（权重10）、点号编号（权重8，后跟大写字母时增强到10）、年份边界（权重3）和arXiv标识符边界（权重5）。相邻10个字符内的边界信号被合并（保留较高权重），最终需要至少2个有效边界（权重$\geq$3）才执行分割。

此外，系统对分割后的片段进行分类（arXiv尾部、DOI尾部、出版商名称、年份尾部、URL尾部、OCR碎片等），将短片段合并到前一条引用记录中，避免信息丢失。

\subsection{三级匹配引擎}

三级匹配引擎是引用图构建的核心算法，其设计目标是高效、准确地将引用记录中的论文标题与arXiv元数据中的论文进行匹配。

\textbf{Level 0：arXiv标识符直接匹配。}如果引用记录中包含arXiv标识符（格式为\texttt{arXiv:NNNN.NNNNN}或\texttt{arxiv.org/abs/NNNN.NNNNN}），系统直接在标识符映射表中进行O(1)哈希查找。这是最可靠的匹配方式，因为arXiv标识符具有全局唯一性。

\textbf{Level 1：精确标题匹配。}将引用记录中提取的论文标题进行归一化处理（转小写、去除标点、统一空白字符），然后在标题-标识符映射表中进行精确查找。归一化处理消除了大小写和标点差异对匹配的影响。

\textbf{Level 2：模糊标题匹配。}当精确匹配失败时，系统使用基于倒排索引\cite{zobel2006inverted, manning2008ir}的模糊匹配算法。该算法分为三个阶段：

（a）\textit{候选生成。}将归一化后的引用标题分词，去除停用词（31个常见英文停用词）和长度小于3的词，保留"显著词"。对每个显著词，通过倒排索引查找包含该词的所有元数据标题，统计每个候选标题与查询标题的显著词重叠数。仅保留重叠数$\geq$2的候选标题。

（b）\textit{候选过滤。}对候选标题进行长度过滤——如果候选标题的归一化长度与查询标题相差超过50\%，则丢弃该候选。剩余候选按重叠数降序排列，仅评估前500个。

（c）\textit{匹配验证。}对每个候选标题，依次尝试两种匹配策略：

\hspace{1em}\textbullet~\textit{双向子串匹配}：检查两个标题是否存在子串包含关系（长度比率$>$0.75），并验证关键词重叠率$\geq$0.35。

\hspace{1em}\textbullet~\textit{SequenceMatcher模糊匹配}\cite{ratcliff1988gestalt}：使用Python \texttt{difflib.SequenceMatcher}计算两个标题字符串的相似度比率。设置匹配阈值为0.90（即两个标题至少有90\%的字符匹配），并要求关键词重叠率$\geq$0.30。当相似度达到0.95时，直接返回匹配结果（早期退出）。

所有匹配结果都经过arXiv标识符交叉验证：当通过标题匹配到一篇论文时，如果引用记录中也包含arXiv标识符，系统验证匹配论文的标识符与引用记录中的标识符是否一致（关键词重叠率$\geq$0.40），防止误匹配。

以cs.SI领域为例，系统从11,929篇论文中提取了21,638条原始引用边，经过三级匹配引擎处理后，成功匹配4,483条域内引用边，匹配率为20.7\%。未匹配的引用主要指向域外论文（其他领域的论文），这一匹配率在学术引用图的构建中是合理的。

\subsection{并行化与零拷贝优化}

引用图构建需要访问约299万篇论文的元数据索引（标题-标识符映射表、标识符-标题映射表、倒排索引等），这些索引数据结构的总内存占用约为数GB。如果使用Python的\texttt{multiprocessing.Pool}在多个工作进程间共享这些数据，传统的序列化-反序列化方式将产生巨大的开销。

系统采用fork-based Copy-on-Write（COW）零拷贝优化来解决这个问题。在创建进程池之前，主进程加载全部元数据并构建索引数据结构，将其存储为模块级别的全局变量。当\texttt{multiprocessing.Pool}使用\texttt{fork()}系统调用创建子进程时，子进程通过操作系统的Copy-on-Write机制继承父进程的内存空间。由于子进程只读取这些数据而不修改它们，父进程和所有子进程共享同一份物理内存页，实现了真正的零拷贝共享。这种优化避免了GB级别的数据序列化开销，显著提升了并行引用图构建的效率。

\section{并行化与性能优化}

本节总结系统中的并行策略及其性能优化效果。

\subsection{多GPU并行提取}

多模态提取是整个系统中计算开销最大的环节。系统设计了\texttt{GPUExtractionPool}实现多GPU并行批量提取（详见3.4.2节），将待处理的论文列表按轮转方式分配到多个GPU上，每个GPU由独立线程驱动，并通过批量处理策略将模型加载开销从每篇30秒降至每批30秒。

\subsection{多进程并行清洗}

数据清洗阶段采用\texttt{multiprocessing.Pool}创建指定数量的CPU工作进程（默认32个），每个进程独立执行完整的六步清洗流水线。由于清洗过程是纯CPU操作且无共享状态，多进程并行可以实现近线性的加速比。

\subsection{COW零拷贝优化}

引用图构建需要访问约299万篇论文的元数据索引（标题-标识符映射表、倒排索引等），总内存占用约为数GB。系统采用fork-based Copy-on-Write（COW）零拷贝优化来避免GB级别数据的序列化开销。在创建进程池之前，主进程加载全部元数据并构建索引数据结构，将其存储为模块级别的全局变量。当\texttt{multiprocessing.Pool}使用\texttt{fork()}系统调用创建子进程时，子进程通过操作系统的Copy-on-Write机制继承父进程的内存空间。由于子进程只读取这些数据而不修改它们，父进程和所有子进程共享同一份物理内存页，实现了真正的零拷贝共享。

\subsection{性能对比}

表\ref{tab:performance}总结了不同配置下的处理性能。以cs.SI领域（11,929篇论文）为例，从元数据过滤到最终数据集输出的完整流程，在8 GPU并行提取和32 CPU并行清洗的配置下，相比串行模式实现了约7倍的总处理时间缩减。

\begin{table}[htbp]
  \centering
  \caption{不同配置下的处理性能（cs.SI, 11,929篇论文）}
  \label{tab:performance}
  \begin{tabular}{lccc}
    \toprule
    配置 & 模型加载 & 提取并行度 & 估算总时间 \\
    \midrule
    串行（1 GPU, 1 CPU） & 每篇30s & 1 & $\sim$99h \\
    多GPU批量（8 GPU） & 每批30s & 8 & $\sim$14h \\
    \bottomrule
  \end{tabular}
\end{table}

\section{MCP工具管理层}

MCP工具管理层是系统三层架构的基础层，负责管理和调度所有底层工具。工具管理器采用插件式注册模式，支持四类共13个工具接口。

\textbf{论文下载工具组：}包括单篇下载工具（支持重试和指数退避）、批量下载工具（使用线程池并行下载，最大10个线程）和DOI下载工具。下载工具在下载前移除所有代理环境变量，避免代理配置导致的连接问题。

\textbf{多模态提取工具组：}包括单篇提取工具、批量提取工具以及文本/图片/表格提取工具。核心提取工具封装了PDF-Extract-Kit的子进程调用，设置30分钟超时防止处理卡死。

\textbf{数据清洗工具组：}包括文本清洗、数据去重和数据验证三个工具接口。

\textbf{数据集工具组：}包括属性图格式化和数据集保存两个工具接口。

工具管理器的核心接口为\texttt{call\_tool(name, **kwargs)}，智能体通过工具名称和参数调用对应的工具函数。这种设计使得新增工具只需实现函数签名并注册到管理器中，无需修改智能体或编排器的代码，实现了工具层与智能体层的完全解耦。这种设计借鉴了AutoGen\cite{wu2023autogen}等框架的工具管理思想，通过标准化的工具接口实现工具的即插即用——其设计理念与模型上下文协议（Model Context Protocol, MCP）类似。

\section{本章小结}

本章详细阐述了基于多智能体系统（MAS）的自动化构建系统的三层架构设计。系统采用编排器-智能体-工具三层多智能体架构，包含五个专用智能体，通过编排器的中央协调实现了从arXiv元数据到标准化多模态属性图数据集的端到端自动化流程。多智能体系统架构的引入是本系统设计的核心决策——学术Benchmark数据集的构建涉及多个异构处理阶段，多智能体系统通过职责分离使每个智能体专注于单一任务，通过松耦合的接口设计实现灵活编排，通过独立的性能优化（如多GPU并行提取、多进程并行清洗、COW零拷贝）使每个智能体在其特定任务上达到最优性能。这种架构模式使得系统具备良好的可扩展性——未来新增处理阶段（如表格结构提取、LLM辅助引用匹配等）只需实现新的智能体并注册到编排器中。

系统的核心技术创新包括：（1）基于DocLayout-YOLO、UniMERNet和PaddleOCR的多模态提取引擎，实现了论文PDF中文本、图像和公式的并行提取；（2）双源引用提取策略，通过PDF直接解析和Markdown提取两种互补数据源，将引用信息的完整性提升了209\%；（3）三级匹配引擎，结合arXiv标识符直接匹配、精确标题匹配和基于倒排索引的模糊匹配，实现了高效的引用-论文匹配；（4）多GPU并行提取和多进程并行清洗，通过批量处理和进程池策略实现了高效的数据处理；（5）fork-based COW零拷贝优化，避免了GB级别元数据索引的序列化开销。

% 第四章 数据集构建与分析
%%
% 第四章 数据集构建与分析（~8页）← 重点章
%%

\chapter{数据集构建与分析}

本章介绍AMMAG数据集的构建过程和统计分析。首先阐述数据采集策略和领域选择依据，然后对数据集的拓扑结构、模态分布和数据质量进行全面的统计分析，最后将AMMAG与现有代表性数据集进行系统对比，展示其独特价值。

\section{数据采集策略}

\subsection{领域选择依据}

AMMAG数据集旨在覆盖多个学科大类，以验证多模态属性图学习方法在不同学术领域的适用性。数据集选择了四个arXiv子领域，横跨计算机科学、数学和物理学三个学科大类。表\ref{tab:domain_design}展示了各子领域的设计和选择理由。

\begin{table}[htbp]
  \centering
  \caption{AMMAG数据集领域设计}
  \label{tab:domain_design}
  \begin{tabular}{llclp{4.5cm}}
    \toprule
    数据集 & 子领域 & 学科大类 & 论文数 & 选择理由 \\
    \midrule
    AMMAG-SI & cs.SI（社交信息学） & CS & 11,929 & 偏社交网络研究，图文并茂 \\
    AMMAG-LO & cs.LO（计算逻辑） & CS & 11,012 & 偏计算理论，公式密度极高 \\
    AMMAG-MLO & math.LO（数理逻辑） & Math & 9,555 & 数学领域，公式密度最高 \\
    AMMAG-SOC & physics.soc-ph（社会物理） & Physics & 13,176 & 物理领域，实验图表丰富 \\
    \bottomrule
  \end{tabular}
\end{table}

领域选择遵循以下三个原则：

\textbf{（1）学科大类覆盖。}四个子领域覆盖了计算机科学、数学和物理学三个学科大类，能够验证多模态图学习方法在不同学科背景下的泛化能力。不同学科的论文在写作风格、模态组成和引用模式上存在显著差异，这种差异为模态消融实验提供了天然的实验条件。

\textbf{（2）模态分布差异。}四个子领域在模态分布上具有鲜明的差异性。cs.SI领域的论文以社交网络分析为主，平均每篇论文包含12.41张图片和180.24条公式，图文比例相对均衡。cs.LO领域的论文偏重计算理论，平均公式密度高达874.70条/篇，而图片仅有6.58张/篇，公式远多于图片。math.LO领域作为纯数学领域，公式密度最高（1,136.50条/篇），图片最少（1.83张/篇）。physics.soc-ph领域的论文涉及物理实验和模拟，图片较多（10.81张/篇），文本相对较短。这种模态分布的差异使得AMMAG能够系统地评估不同模态在不同学术背景下的贡献。

\textbf{（3）领域内差异验证。}在计算机科学内部选择了两个差异显著的子领域：cs.SI偏应用（社交网络分析、推荐系统），cs.LO偏理论（计算复杂性、形式逻辑）。这种设计使得跨域迁移实验（第五章5.7节）能够在同属CS但子领域差异显著的场景下进行，更精细地评估模型的迁移能力。

\subsection{采集流程}

数据采集流程基于第三章所述的自动化构建系统实现，具体步骤如下：

\textbf{（1）元数据过滤。}从arXiv元数据集（约299万篇论文）中，按主领域标签分别过滤出cs.SI、cs.LO、math.LO和physics.soc-ph四个子领域的论文，按更新日期降序排序后各取前12,000篇作为候选集。

\textbf{（2）PDF批量下载。}使用多线程下载工具（最大10个并发线程）批量下载候选论文的PDF文件。四个子领域的下载成功率均在99\%以上：cs.SI为99.6\%（11,951/12,000），cs.LO为99.3\%（11,017/11,100），math.LO为98.9\%（9,555/9,660），physics.soc-ph为99.8\%（13,176/13,200）。

\textbf{（3）多模态提取。}使用基于DocLayout-YOLO\cite{zhao2024doclayoutyolo}的多模态提取引擎，从每篇PDF中提取文本（Markdown格式）、图片（PNG格式）和公式（LaTeX格式）。提取过程采用多GPU并行模式（6个GPU，批次大小50），利用批量处理策略将模型加载开销从每篇30秒降至每批30秒。

\textbf{（4）数据清洗。}对提取的Markdown文本执行六步清洗流水线（详见第三章3.4节），修复PDF转换引入的格式问题，提取和格式化参考文献列表。清洗过程采用32个CPU工作进程并行执行。

\textbf{（5）引用图构建。}使用ReferenceGraphBuilderAgent的双源提取和三级匹配引擎，从每篇论文中提取引用信息并与arXiv元数据进行匹配，构建域内引用网络图。

\textbf{（6）标准化格式化。}将清洗后的多模态数据、引用边信息和论文元数据整合为标准化的多模态属性图数据集，输出为PyTorch Geometric兼容的格式。

\section{数据集统计分析}

\subsection{基本统计}

表\ref{tab:statistics}展示了AMMAG数据集四个子领域的基本统计信息。

\begin{table}[htbp]
  \centering
  \caption{AMMAG数据集统计}
  \label{tab:statistics}
  \begin{tabular}{lrrrrr}
    \toprule
    统计项 & cs.SI & cs.LO & math.LO & physics.soc-ph & 合计 \\
    \midrule
    论文数 & 11,929 & 11,012 & 9,555 & 13,176 & 45,672 \\
    引用边数（域内） & 4,483 & 8,174 & 12,544 & 30,325 & 55,526 \\
    引用边数（含跨域） & --- & --- & --- & --- & 59,364 \\
    平均度 & 0.75 & 1.49 & 2.63 & 4.60 & 2.60 \\
    \midrule
    平均图片数/篇 & 12.41 & 6.58 & 1.83 & 10.81 & 7.91 \\
    平均公式数/篇 & 180.24 & 874.70 & 1,136.50 & 188.70 & 594.79 \\
    平均文本长度(字符) & 61,361 & 100,214 & 87,735 & 53,996 & 75,826 \\
    \midrule
    含图片论文比例 & 98.5\% & 81.6\% & 38.2\% & 97.3\% & 81.5\% \\
    含公式论文比例 & 99.4\% & 99.8\% & 100.0\% & 99.6\% & 99.7\% \\
    \bottomrule
  \end{tabular}
\end{table}

从表\ref{tab:statistics}可以观察到以下关键特征。图\ref{fig:dataset_overview}直观展示了AMMAG数据集的引用网络结构和模态分布。

\begin{figure}[htbp]
  \centering
  \includegraphics[width=0.92\textwidth]{images/dataset_overview.png}
  \caption{AMMAG数据集概览：引用网络结构与模态分布}
  \label{fig:dataset_overview}
\end{figure}

\textbf{（1）图结构特征。}AMMAG合并引用图包含45,672个节点和59,364条引用边（含跨领域引用3,838条），平均度为2.60。四个子领域的图结构密度差异显著：physics.soc-ph最稠密（平均度4.60），cs.SI最稀疏（平均度0.75）。这种差异反映了不同学科的引用习惯——物理学论文倾向于引用更多相关领域的工作，而计算机科学的某些子领域（如社交信息学）的引用范围相对集中。

\textbf{（2）模态分布差异。}四个子领域在模态分布上呈现出鲜明的差异性。在图片方面，cs.SI的论文平均包含12.41张图片（主要为网络架构图和实验结果可视化），而math.LO仅包含1.83张图片（数学论文以纯文本和公式为主），两者相差约6.8倍。在公式方面，math.LO的论文平均包含1,136.50条公式，而cs.SI仅包含180.24条，两者相差约6.3倍。在文本长度方面，cs.LO的论文最长（平均100,214字符），physics.soc-ph最短（53,996字符）。

\textbf{（3）模态覆盖率。}公式模态的覆盖率最高，四个子领域中99.7\%的论文至少包含一条公式，这符合学术论文的基本特征——即使是实验性论文也通常包含数学符号和简单公式。图片模态的覆盖率因领域而异：cs.SI和physics.soc-ph的覆盖率超过97\%，而math.LO仅有38.2\%的论文包含图片。

\subsection{与现有数据集的系统对比}

表\ref{tab:benchmark_comparison}将AMMAG与现有代表性图基准数据集进行了系统对比。

\begin{table}[htbp]
  \centering
  \caption{AMMAG与现有Benchmark数据集对比}
  \label{tab:benchmark_comparison}
  \small
  \begin{tabular}{p{2.2cm}p{1.2cm}rrcccp{1.2cm}}
    \toprule
    数据集 & 领域 & 节点数 & 边数 & 文本 & 图像 & 公式 & 图结构 \\
    \midrule
    Cora\cite{sen2008collective} & 学术 & 2,708 & 5,429 & BoW & \texttimes & \texttimes & 引用网络 \\
    CiteSeer\cite{giles1998citeseer} & 学术 & 3,327 & 4,732 & BoW & \texttimes & \texttimes & 引用网络 \\
    PubMed\cite{sen2008collective} & 学术 & 19,717 & 44,338 & TF-IDF & \texttimes & \texttimes & 引用网络 \\
    ogbn-arxiv\cite{hu2020ogb} & 学术 & 169,343 & 1,166,243 & 128d向量 & \texttimes & \texttimes & 引用网络 \\
    CS-TAG\cite{yan2023cstag} & 学术 & 18,000 & 82,000 & 原始文本 & \texttimes & \texttimes & 合作网络 \\
    DocGenome\cite{xia2024docgenome} & 科学 & 500,000 & --- & 全文 & 原始 & LaTeX & \texttimes \\
    \midrule
    MAGB-Movies\cite{yan2025magb} & 电商 & 16,238 & 218,572 & 原始 & 原始 & \texttimes & 共购网络 \\
    MAGB-Reddit\cite{yan2025magb} & 社交 & 15,000 & 566,632 & 原始 & 原始 & \texttimes & 互动网络 \\
    OpenMAG\cite{wan2026openmag} & 多领域 & 多规模 & --- & 原始 & 原始 & \texttimes & 部分 \\
    MM-GRAPH\cite{zhu2025mmgraph} & 多领域 & 多规模 & --- & 原始 & 原始 & \texttimes & 部分 \\
    \midrule
    \textbf{AMMAG} & \textbf{学术} & \textbf{45,672} & \textbf{59,364} & \textbf{全文MD} & \textbf{原始} & \textbf{LaTeX} & \textbf{引用网络} \\
    \bottomrule
  \end{tabular}
\end{table}

从表\ref{tab:benchmark_comparison}的对比中可以清晰地看到AMMAG的独特定位：

\textbf{与经典学术图数据集的对比。}Cora\cite{sen2008collective}、CiteSeer\cite{giles1998citeseer}和PubMed\cite{sen2008collective}作为最早被广泛使用的学术图数据集，其节点属性仅为词袋或TF-IDF表示，完全丢失了原始文本的语义信息，更不包含图像和公式等多模态信息。AMMAG提供了原始的全文Markdown文本、论文图片的PNG文件和公式的LaTeX源码，保留了完整的原始多模态信息。

\textbf{与OGB的对比。}ogbn-arxiv\cite{hu2020ogb}虽然规模更大（169K节点），但其节点特征为128维的skip-gram向量，是将论文标题和摘要通过skip-gram模型编码得到的低维表示。这种压缩表示丢失了大量语义信息——不同领域的论文如果标题中的关键词相似，可能产生接近的特征向量，无法体现论文的完整语义。AMMAG保留了原始多模态数据，研究者可以根据需要选择不同的编码策略（如SBERT\cite{reimers2019sbert}、SciBERT\cite{beltagy2019scibert}等），灵活性更高。

\textbf{与MAG基准的对比。}MAGB\cite{yan2025magb}、OpenMAG\cite{wan2026openmag}和MM-GRAPH\cite{zhu2025mmgraph}虽然提供了多模态节点属性，但仅限于文本和图像两种模态。AMMAG首次将公式模态引入MAG基准。公式在学术论文中承载了大量的结构化数学信息，是数学、物理和理论计算机科学论文最核心的内容之一。缺少公式模态使得现有MAG基准无法评估公式在学术图学习任务中的贡献。

\subsection{模态分布分析}

为了更直观地展示四个子领域的模态分布差异，本节从多个角度进行深入分析。

\textbf{图片数量分布。}四个子领域的图片数量分布均呈现长尾分布特征，但分布形态存在显著差异。cs.SI和physics.soc-ph领域的图片数量分布相对分散，部分论文包含超过50张图片（主要为实验结果可视化），而中位数约为8-10张。math.LO领域约53.7\%的论文不包含任何图片，有图片的论文中位数仅为1-2张，且图片多为定理证明中的辅助示意图。cs.LO领域介于两者之间，图片数量集中在2-10张的范围内。

\textbf{公式数量分布。}公式数量的跨领域差异更为极端。math.LO领域的公式数量分布具有最长的尾部——部分论文包含超过5,000条公式（主要为形式化证明的论文），平均值为1,136.50条。cs.LO领域的公式密度也很高（平均874.70条/篇），因为计算逻辑领域的论文大量使用形式化规约和证明。相比之下，physics.soc-ph和cs.SI的公式密度较低（分别平均188.70和180.24条/篇），这些领域的论文更多使用叙述性文字和实验图表来传达研究结果。

\textbf{度分布特征。}四个子领域的度分布均呈现幂律分布特征——少数"hub"论文拥有大量的引用，而大部分论文的引用次数较少。physics.soc-ph的度分布尾部最长，存在度数超过100的hub论文；cs.SI的度分布最为集中，大部分论文的度数在0-3之间。图\ref{fig:degree_dist}直观展示了四个子领域的度分布差异。

\begin{figure}[htbp]
  \centering
  \includegraphics[width=0.95\textwidth]{images/degree_distribution.png}
  \caption{四个子领域的引用网络度分布（对数坐标）}
  \label{fig:degree_dist}
\end{figure}

图\ref{fig:modality_dist}以分组柱状图的形式直观对比了四个子领域在图片数量、公式数量和文本长度三个维度上的差异。

\begin{figure}[htbp]
  \centering
  \includegraphics[width=0.95\textwidth]{images/modality_distribution.png}
  \caption{四个子领域的模态分布对比}
  \label{fig:modality_dist}
\end{figure}

\section{数据质量评估}

为了评估AMMAG数据集的数据质量，本节从多个维度进行量化分析。

\subsection{多模态提取质量}

表\ref{tab:quality}展示了数据质量的评估结果。

\begin{table}[htbp]
  \centering
  \caption{数据质量评估}
  \label{tab:quality}
  \begin{tabular}{lll}
    \toprule
    评估维度 & 评估指标 & 结果 \\
    \midrule
    文本完整性 & 提取后文本长度与PDF页数的相关系数 & $\sim$0.85 \\
    图片提取率 & 成功提取图片的论文比例 & 81.5\% (37,210/45,672) \\
    公式提取率 & 成功提取公式的论文比例 & 99.7\% (45,531/45,672) \\
    引用边质量 & 域内匹配率 & 20.7\% (cs.SI) \\
    数据清洗效果 & 重复段落去除率 & $>$90\% \\
    \bottomrule
  \end{tabular}
\end{table}

\textbf{文本完整性。}提取后Markdown文本的长度（字符数）与原始PDF的页数之间的皮尔逊相关系数约为0.85，表明文本提取的完整性与论文的实际内容量高度相关。较长的论文产生较长的提取文本，这说明提取过程没有系统性地遗漏大段文本内容。相关系数未达到1.0的原因包括：部分论文包含大量图片（图片区域不产生文本）、参考文献列表的提取完整度因格式差异而不同、以及双栏布局论文的文本合并可能产生少量信息损失。

\textbf{图片提取率。}81.5\%的论文至少包含一张成功提取的图片。图片提取率因领域而异：cs.SI为98.5\%（该领域论文通常包含网络架构图和实验可视化），math.LO仅为38.2\%（纯数学论文通常不包含图片）。对于不含图片的论文，图像编码器输出768维零向量作为占位符，确保这些论文仍然可以参与多模态图学习实验。

\textbf{公式提取率。}99.7\%的论文至少包含一条成功提取的公式，这一高覆盖率反映了公式在学术论文中的普遍性——即使是非数学领域的论文也通常包含符号定义、目标函数或评价指标公式等数学内容。公式提取使用UniMERNet\cite{wang2024unimernet}将公式图像识别为LaTeX代码，其识别精度在标准数学表达式上表现优异。

\textbf{引用边质量。}以cs.SI领域为例，系统从11,929篇论文中提取了21,638条原始引用关系，经过三级匹配引擎处理后成功匹配4,483条域内引用边，匹配率为20.7\%。未匹配的引用主要指向域外论文（即被引用论文不在当前数据集中）。这一较低的匹配率限制了图的密度（合并图平均度仅2.60），直接影响了GNN图结构增益的体现——第五章实验中GCN与MLP仅差0.26pp，可能部分归因于图的稀疏性。跨领域引用的高丢失率是学术引用网络的固有特征——一篇cs.SI论文可能引用机器学习（cs.LG）、信息检索（cs.IR）等多个领域的工作。

\subsection{标签体系}

AMMAG数据集提供两级标签体系：

\textbf{领域标签（4类）。}每篇论文根据其所属的arXiv主领域被赋予一个领域标签：cs.SI（0）、cs.LO（1）、math.LO（2）、physics.soc-ph（3）。领域标签用于节点分类任务（第五章5.2节）和模态消融实验（第五章5.3节）。四个领域的论文数量分布相对均衡（9,555-13,176篇），避免了严重的类别不平衡问题。

\textbf{子分类标签（6-7类/领域）。}每篇论文根据其在arXiv中的交叉类别（cross-list categories）被赋予一个子分类标签。子分类标签的分配规则为：从论文的类别列表中选择第一个非主领域的、且在该领域的子分类映射表中存在的类别；如果没有匹配的交叉类别，则归入"other"类别。各领域的子分类数量和分布如下：cs.SI有6个子分类（cs.SI、cs.LG、cs.CY、cs.CL、cs.AI、other），cs.LO有7个子分类，math.LO有7个子分类，physics.soc-ph有6个子分类。子分类标签用于更细粒度的分类任务（第五章5.7节）和跨域迁移实验（第五章5.6节）。

\section{特征编码与数据集格式}

\subsection{统一编码方案}

AMMAG数据集的所有模态统一使用CLIP\cite{radford2021clip}（Contrastive Language-Image Pre-training）的ViT-Large-Patch14模型进行编码，输出维度均为768维。该统一编码方案确保了文本、图像和公式三种模态在同一语义空间中的表示，为跨模态融合和检索提供了基础。

\textbf{文本编码。}使用CLIP的文本编码器将论文的标题与摘要编码为768维向量。对于长文档，编码器将文本按段落分割为多个片段分别编码，最后取均值池化作为整篇论文的文本表示。CLIP文本编码器基于Transformer架构，在大规模图像-文本对上进行了对比学习预训练，能够捕获丰富的语义信息。

\textbf{图像编码。}使用CLIP的视觉编码器（ViT-Large-Patch14）将论文中的图片编码为768维向量。对于包含多张图片的论文，对每张图片分别编码后取均值池化。不含图片的论文返回768维零向量作为占位符。CLIP视觉编码器采用Vision Transformer架构，将图像分割为$14\times14$的patch后进行自注意力编码。

\textbf{公式编码。}使用与文本编码器相同的CLIP文本编码器将论文中的LaTeX公式编码为768维向量。编码器从Markdown文本中提取块级公式和行内公式，过滤短公式后取前200条，分别编码后取均值池化。将LaTeX公式作为文本序列用CLIP编码，利用了CLIP文本编码器的通用序列建模能力来处理数学符号。

\textbf{统一编码的优势。}与使用不同编码器分别编码各模态的方案相比，统一编码方案具有两个关键优势：（1）跨模态语义对齐——CLIP的对比学习预训练使得文本和图像在共享的768维空间中具有语义对应关系，跨模态检索MRR（11.45\%）是独立编码器（0.17\%）的67倍（第五章5.8.2节详细验证）；（2）维度一致性——所有模态输出均为768维，避免了不同维度（如SBERT 384维与CLIP 512维）拼接时的模态失衡问题。

\subsection{多模态融合策略}

AMMAG数据集支持7种模态组合的节点特征表示。当使用多模态组合时，各模态的768维向量直接拼接后，通过一个可学习的线性投影层映射至统一的768维空间：

\begin{equation}
\mathbf{x}_i = \mathbf{W}_{proj} \cdot [\mathbf{x}_i^{text}; \mathbf{x}_i^{image}; \mathbf{x}_i^{formula}] + \mathbf{b}_{proj}
\end{equation}

其中$\mathbf{W}_{proj} \in \mathbb{R}^{768 \times d_{in}}$，$d_{in}$为拼接后的维度（单模态为768，双模态为1536，三模态为2304）。该投影层在下游GNN模型训练过程中端到端学习，确保不同维度的模态组合能够公平地与768维单模态进行对比。七种模态组合及其特征维度为：Text（768）、Image（768）、Formula（768）、T+I（1536→768）、T+F（1536→768）、I+F（1536→768）、T+I+F（2304→768）。

\subsection{标准化数据格式}

AMMAG数据集以PyTorch Geometric（PyG）\cite{fey2019pyg}兼容的格式存储，每个领域的输出目录包含以下文件：

\begin{itemize}
  \item \texttt{graph\_data.pt}：PyG Data对象，包含\texttt{edge\_index}（引用边的稀疏张量）、\texttt{labels}（领域标签）、\texttt{subcat\_labels}（子分类标签）和\texttt{paper\_ids}（论文ID列表）。
  \item \texttt{text\_embeddings.pt}：文本模态的编码结果，维度为$N \times 768$，与\texttt{paper\_ids}对齐。
  \item \texttt{image\_embeddings.pt}：图像模态的编码结果，维度为$N \times 768$。
  \item \texttt{formula\_embeddings.pt}：公式模态的编码结果，维度为$N \times 768$。
  \item \texttt{statistics.json}：数据集的统计元信息。
\end{itemize}

数据集的使用接口设计简洁，研究者可以通过以下方式加载数据：

\begin{lstlisting}[language=Python, caption={AMMAG数据集加载示例}]
from benchmark.data.dataset import AMMAGDataset

# 加载单个领域
dataset = AMMAGDataset(
    root='./data', domains=['cs_si'],
    modality='text+formula', label_type='domain'
)

# 加载合并的4领域图
dataset = AMMAGDataset(
    root='./data',
    domains=['cs_si', 'cs_lo', 'math_lo', 'physics_socph'],
    modality='text+image+formula'
)

# 获取数据划分
data = dataset._load_data()
train_mask, val_mask, test_mask = dataset.get_split_masks('supervised')
\end{lstlisting}

\texttt{AMMAGDataset}类支持多种数据划分方式：（1）全监督划分——随机打乱后按60\%/20\%/20\%比例划分训练集、验证集和测试集；（2）Few-shot划分——每类随机采样$k$个样本作为训练集（$k \in \{1, 3, 5, 10, 20\}$），20\%作为验证集，其余作为测试集；（3）跨域迁移划分——以非目标领域的所有节点作为训练集和验证集（80\%/20\%），目标领域的节点作为测试集；（4）域内子图划分——提取单个领域的子图进行独立评估。

\section{本章小结}

本章介绍了AMMAG数据集的构建过程和统计分析。AMMAG覆盖计算机科学、数学和物理学三个学科大类共45,672篇论文，包含文本、图像和公式三种模态属性，并通过引用网络连接。

数据集的核心特征包括：（1）四个子领域在模态分布上差异显著——math.LO的公式密度是cs.SI的6.3倍，cs.SI的图片数量是math.LO的6.8倍，这种差异为模态消融实验提供了理想的实验条件；（2）图结构密度差异明显——physics.soc-ph的平均度（4.60）是cs.SI（0.75）的6.1倍，反映了不同学科的引用习惯差异；（3）公式模态覆盖率极高（99.7\%），验证了公式作为学术论文核心模态的地位。

在特征编码方面，AMMAG采用CLIP-ViT-Large-Patch14统一编码方案，将文本、图像和公式三种模态统一映射至768维语义空间，并使用可学习线性投影层将多模态拼接特征投影至统一维度。该方案相比独立编码器方案具有显著的跨模态检索优势（MRR提升67倍），为多模态融合提供了高质量的语义对齐基础。

与现有数据集的系统对比表明，AMMAG是首个同时满足学术论文领域、三模态节点属性（含公式）、引用网络图结构和自动化构建系统四个条件的数据集。AMMAG继承了OGB\cite{hu2020ogb}和CS-TAG\cite{yan2023cstag}的学术领域和图结构优势，借鉴了DocGenome\cite{xia2024docgenome}的多模态信息提取技术，融合了MAGB\cite{yan2025magb}的MAG学习范式，并通过第三章所述的自动化构建系统实现了端到端的数据集生成。

% 第五章 Benchmark实验与分析
%%
% 第五章 Benchmark实验与分析（~12页）← 最重点章
%%

\chapter{Benchmark实验与分析}

本章在AMMAG数据集上进行全面的Benchmark实验评估。首先介绍实验设置，包括统一编码方案、基线模型和超参数配置。然后从多模态必要性、模型有效性、鲁棒性、Few-shot学习、链接预测和跨域迁移六个维度展开实验分析，全面评估AMMAG数据集上多模态图学习的性能与特性。

本章评估的实验围绕以下核心任务展开：

\textbf{全监督节点分类。}给定多模态属性图$\mathcal{G}=(\mathcal{V}, \mathcal{E}, \mathbf{X})$，其中$\mathcal{V}$为论文节点集合，$\mathcal{E}$为引用边集合，$\mathbf{X} \in \mathbb{R}^{|\mathcal{V}| \times d}$为节点多模态特征矩阵，目标是学习映射$f: \mathcal{V} \rightarrow \mathcal{Y}$，将每个论文节点分类到预定义的标签集合$\mathcal{Y}$中（领域4类或子分类6-7类）。

\textbf{链接预测。}给定图$\mathcal{G}$中已知的引用边集合$\mathcal{E}_{train}$，目标是学习评分函数$s: \mathcal{V} \times \mathcal{V} \rightarrow \mathbb{R}$，使得存在的引用边的得分高于不存在的边，使用MRR和Hits@K作为评估指标。

\textbf{Few-shot节点分类。}在极端少样本条件下进行节点分类，每类仅使用$k$个标注样本（$k \in \{1, 3, 5, 10, 20\}$）作为训练集，评估模型在低数据条件下的泛化能力。

\textbf{跨域迁移。}给定源域$\mathcal{D}_s$和目标域$\mathcal{D}_t$，先在$\mathcal{D}_s$子图上预训练GNN编码器，再将编码器权重复制到目标域模型中进行微调，评估不同学术领域间的知识可迁移性。

\section{实验设置}

\subsection{统一编码方案}

所有模态统一使用CLIP\cite{radford2021clip} ViT-Large-Patch14编码至768维向量空间，多模态组合通过可学习线性投影层映射至统一768维。表\ref{tab:encoders}展示了编码器配置。

\begin{table}[htbp]
  \centering
  \caption{特征编码器配置}
  \label{tab:encoders}
  \begin{tabular}{llcp{4.5cm}}
    \toprule
    模态 & 编码器 & 输出维度 & 编码策略 \\
    \midrule
    文本 (Text) & CLIP\cite{radford2021clip} ViT-Large & 768 & 标题+摘要分段编码$\to$均值池化 \\
    图像 (Image) & CLIP ViT-Large & 768 & 多图编码$\to$均值池化（无图论文用零向量） \\
    公式 (Formula) & CLIP ViT-Large & 768 & LaTeX公式提取（最多200条）$\to$编码$\to$均值池化 \\
    多模态融合 & 拼接+投影 & 768 & 各模态拼接后通过Linear($d_{in}$, 768)投影 \\
    \bottomrule
  \end{tabular}
\end{table}

该统一编码方案的核心优势在于：CLIP的对比学习预训练使文本和图像在同一768维空间中具有语义对应关系，跨模态检索MRR（11.45\%）是独立编码器（SBERT+CLIP-Base，0.17\%）的67倍（详见5.8.2节）。所有模态输出维度一致（768维），避免了不同维度拼接时的模态失衡问题。

\subsection{基线模型}

本文评估了13种基线模型（另有GT仅参与Few-shot实验），涵盖标准GNN和多模态GNN两大类：

\textbf{标准模型（8种，含MLP和GT）：}

\textbf{MLP}：多层感知机作为无图结构的基线模型，仅使用节点特征进行分类。MLP的性能反映特征本身的信息量。

\textbf{GCN}\cite{kipf2017gcn}：图卷积网络，通过归一化邻接矩阵实现谱域近似图卷积。

\textbf{GAT}\cite{velickovic2018gat}：图注意力网络，使用多头注意力机制（4头）自适应学习邻居权重。

\textbf{GraphSAGE}\cite{hamilton2017graphsage}：采样聚合网络，通过采样和聚合策略支持归纳式学习。

\textbf{GIN}\cite{xu2019gin}：图同构网络，基于Weisfeiler-Lehman测试理论设计求和聚合策略。

\textbf{GCNII}：深层图卷积网络，引入初始残差连接与恒等映射，有效缓解深层GNN的过平滑问题。

\textbf{GATv2}：动态注意力图网络，改进GAT的静态注意力机制为动态注意力。

\textbf{GT}：Graph Transformer，使用PyG\cite{fey2019pyg}的TransformerConv实现，配置4头注意力。

\textbf{多模态GNN模型（6种）：}

\textbf{DGF}：解耦图滤波网络，将特征传播与变换解耦。

\textbf{MMGCN}：多模态图卷积网络，对每种模态分别进行图聚合后融合。

\textbf{MGAT}：多模态图注意力网络，在图注意力机制中融入模态感知能力。

\textbf{LGMRec}：线性图多模态推荐网络，采用线性图传播与模态解耦策略。

\textbf{DMGC}：动态多模态图组合网络，动态学习模态组合策略。

\textbf{MHGAT}：多头异质图注意力网络，使用多头机制处理多模态异质信息。

所有模型使用相同网络结构：2层GNN，隐藏维度256，Dropout率0.5。多模态输入通过可学习Linear($d_{in}$, 768)投影层映射至统一维度。

\subsection{超参数与评估设置}

表\ref{tab:hyperparams}展示了实验的超参数设置。

\begin{table}[htbp]
  \centering
  \caption{超参数设置}
  \label{tab:hyperparams}
  \begin{tabular}{ll}
    \toprule
    参数 & 值 \\
    \midrule
    隐藏维度 & 256 \\
    GNN层数 & 2 \\
    Dropout & 0.5 \\
    学习率 & 0.001 \\
    权重衰减 & $5 \times 10^{-4}$ \\
    优化器 & Adam \\
    训练轮数 & 500（early stopping, patience=25） \\
    投影维度 & 768 \\
    监督学习重复次数 & 5次（显著性检验10次） \\
    Few-shot重复次数 & 10次 \\
    跨域迁移重复次数 & 3次 \\
    数据划分 & 60\%训练/20\%验证/20\%测试 \\
    \bottomrule
  \end{tabular}
\end{table}

评估指标包括：分类任务使用准确率（Accuracy）和宏平均F1值（Macro-F1），链接预测任务使用MRR和Hits@K（K=1,3,10）。所有实验使用多次独立运行的均值$\pm$标准差报告结果。

实验使用两级标签：（1）\textbf{领域标签}（4类）——cs.SI、cs.LO、math.LO和physics.soc-ph，用于全监督分类和模态消融实验；（2）\textbf{子分类标签}（每域6-7类），用于跨域迁移实验。

\section{多模态消融实验}

本节是本文的核心创新实验。在7种模态组合和4个基线模型（MLP、GCN、GAT、GraphSAGE）上系统评估多模态融合的必要性，并分析不同学术领域的模态偏好差异。

\subsection{七种模态组合消融}

表\ref{tab:nc_7modality}展示了全监督节点分类（4类领域分类）在7种模态组合下的完整结果。

\begin{table}[htbp]
  \centering
  \caption{七种模态组合的节点分类结果（Accuracy \%，5次运行均值$\pm$标准差）}
  \label{tab:nc_7modality}
  \small
  \begin{tabular}{lccccccc}
    \toprule
    模型 & Text & Image & Formula & T+I & T+F & I+F & T+I+F \\
    \midrule
    MLP & 88.86$\pm$0.24 & 79.47$\pm$0.42 & 76.47$\pm$0.21 & 89.94$\pm$0.36 & \textbf{90.55$\pm$0.25} & 85.06$\pm$0.37 & 90.33$\pm$0.27 \\
    GCN & 89.47$\pm$0.29 & 81.25$\pm$0.28 & 78.35$\pm$0.16 & 90.37$\pm$0.16 & \textbf{90.90$\pm$0.11} & 86.29$\pm$0.37 & 90.67$\pm$0.23 \\
    GAT & 89.42$\pm$0.21 & 81.58$\pm$0.46 & 78.68$\pm$0.20 & 90.40$\pm$0.23 & \textbf{90.83$\pm$0.16} & 86.13$\pm$0.34 & 90.55$\pm$0.20 \\
    GraphSAGE & 89.43$\pm$0.22 & 81.31$\pm$0.18 & 78.43$\pm$0.12 & 90.67$\pm$0.16 & \textbf{91.06$\pm$0.22} & 86.22$\pm$0.23 & 90.91$\pm$0.14 \\
    \midrule
    \textbf{均值} & \textbf{89.30} & \textbf{80.90} & \textbf{77.98} & \textbf{90.35} & \textbf{90.84} & \textbf{85.93} & \textbf{90.62} \\
    \bottomrule
  \end{tabular}
\end{table}

从表\ref{tab:nc_7modality}可以得出以下关键发现：

\textbf{发现1：文本是最强单模态。}Text-only在所有模型上均达到约89-90\%的准确率，远超Image-only（约81\%）和Formula-only（约78\%）。CLIP文本编码器已经捕获了论文标题和摘要的丰富语义信息，足以区分4个学科领域。

\textbf{发现2：T+F（文本+公式）是最佳模态组合。}T+F在所有模型上均优于或持平T+I+F，均值达到90.84\%，超过三模态T+I+F（90.62\%）0.22个百分点。公式比图像提供了更多互补增益，尽管Formula-only（77.98\%）低于Image-only（80.90\%）。这说明公式的独特价值在于其与文本的互补性而非独立区分力。

\textbf{发现3：缺失文本是灾难性的。}I+F组合降至约86\%，损失约4个百分点。Image-only和Formula-only分别仅约81\%和78\%，远低于Text-only的约89\%。文本是不可替代的主导模态。

\textbf{发现4：GraphSAGE获益最大。}GraphSAGE在T+F组合下达到91.06\%，为所有配置中的最优结果，比GraphSAGE+Text（89.43\%）提升1.63个百分点。

表\ref{tab:modality_gain}量化了各模态组合相对Text-only的增益。

\begin{table}[htbp]
  \centering
  \caption{模态组合增益分析（相对Text-only，百分点）}
  \label{tab:modality_gain}
  \begin{tabular}{lcp{6cm}}
    \toprule
    模态组合 & 均值增益 (pp) & 说明 \\
    \midrule
    \textbf{T+F} & \textbf{+1.54} & 最大增益：文本+公式组合效果最佳 \\
    T+I+F & +1.32 & 三模态融合，略低于T+F \\
    T+I & +1.05 & 文本+图像组合，增益温和 \\
    I+F & $-$3.37 & 缺失文本导致严重下降 \\
    Image & $-$8.40 & 图像单独信息不足 \\
    Formula & $-$11.32 & 公式单独信息最弱 \\
    \bottomrule
  \end{tabular}
\end{table}

图\ref{fig:modality_bar}以柱状图直观展示了7种模态组合的分类性能对比。

\begin{figure}[htbp]
  \centering
  \includegraphics[width=0.92\textwidth]{images/modality_ablation_bar.png}
  \caption{七种模态组合的节点分类性能对比}
  \label{fig:modality_bar}
\end{figure}

\subsection{领域模态偏好分析}

表\ref{tab:domain_modality}展示了不同模态在四个领域的分类准确率，揭示了各领域的模态偏好差异。

\begin{table}[htbp]
  \centering
  \caption{各领域的模态偏好分析（Accuracy \%，4模型均值）}
  \label{tab:domain_modality}
  \begin{tabular}{lrrrrr}
    \toprule
    模态 & cs.SI & cs.LO & math.LO & physics.soc-ph & Overall \\
    \midrule
    Text-only & 86.67 & 89.36 & 89.78 & 91.29 & 89.30 \\
    Image-only & 82.05 & 71.93 & 82.94 & 85.97 & 80.90 \\
    T+I+F & 88.78 & 90.48 & 91.21 & 91.84 & 90.62 \\
    \midrule
    T+I+F vs Text 增益 & \textbf{+2.11} & +1.12 & \textbf{+1.43} & +0.55 & +1.32 \\
    \bottomrule
  \end{tabular}
\end{table}

四个领域对多模态的响应截然不同：

\textbf{（1）cs.SI从多模态融合获益最大}（+2.11pp）。该领域论文平均含12.41张图片（最多），图片中的网络架构图和系统设计图对分类贡献显著。Image-only在cs.SI上也达到82.05\%，为Image在各领域中最高。

\textbf{（2）math.LO从公式模态获益最大}（+1.43pp）。该领域平均含1,136条公式（全领域最高），公式提供了关键的数学结构信息。值得注意的是，Image-only在math.LO上仅82.94\%，因为该领域仅38.2\%的论文含图片。

\textbf{（3）physics.soc-ph增益最小}（+0.55pp）。该领域Text-only已达91.29\%（最高），接近性能天花板，多模态增益有限。

\textbf{（4）Image-only在cs.LO上表现最弱}（71.93\%）。计算逻辑领域以逻辑推理为主，图片信息量最低（仅81.6\%论文含图，平均6.58张/篇）。

这一领域差异是AMMAG独有的发现。MAGB\cite{yan2025magb}、OpenMAG\cite{wan2026openmag}和MM-GRAPH\cite{zhu2025mmgraph}因缺少公式模态，无法开展此类分析。

\subsection{图结构信息贡献分析}

为量化图结构信息的贡献，使用GraphSAGE+T+I+F，变化每个节点的采样邻居数量$K \in \{0, 2, 4, 6, 8\}$。$K=0$等价于MLP（不使用邻居信息）。

\begin{table}[htbp]
  \centering
  \caption{邻居数量消融实验（GraphSAGE + T+I+F, Accuracy \%）}
  \label{tab:neighbor}
  \begin{tabular}{lrrrrr}
    \toprule
    $K$（邻居数） & 0 & 2 & 4 & 6 & 8 \\
    \midrule
    Accuracy & 90.68$\pm$0.12 & 91.16$\pm$0.07 & \textbf{91.29$\pm$0.07} & 91.18$\pm$0.10 & 91.28$\pm$0.23 \\
    \bottomrule
  \end{tabular}
\end{table}

引入图结构信息即刻带来显著提升：$K=0 \to K=2$提升0.48pp。$K=4$时达到最优，之后趋于饱和，边际收益递减。该数据集平均度仅2.60，$K=4$已覆盖大部分有效邻居。

\subsection{链接预测模态消融}

表\ref{tab:lp_7modality}展示了链接预测在7种模态组合下的结果。

\begin{table}[htbp]
  \centering
  \caption{七种模态组合的链接预测结果（MRR）}
  \label{tab:lp_7modality}
  \begin{tabular}{lccccccc}
    \toprule
    模型 & Text & Image & Formula & T+I & T+F & I+F & T+I+F \\
    \midrule
    MLP & 0.957 & 0.948 & 0.931 & 0.966 & 0.962 & 0.955 & \textbf{0.970} \\
    GCN & 0.974 & 0.973 & 0.953 & 0.983 & 0.978 & 0.977 & \textbf{0.985} \\
    GAT & 0.959 & 0.960 & 0.941 & 0.971 & 0.957 & 0.966 & \textbf{0.974} \\
    \bottomrule
  \end{tabular}
\end{table}

链接预测整体高度饱和（MRR $>$ 0.93），所有模型和模态组合的区分度不如节点分类明显。T+I+F在所有模型上取得最高MRR。链接预测饱和的原因在于：AMMAG引用网络的同质性强（同领域论文高度互联），链接预测退化为领域识别的代理任务。该数据集不适合作为链接预测Benchmark，建议以节点分类作为主要评估指标。

\section{模型有效性实验}

本节评估13种模型在AMMAG上的表现，包括标准GNN和多模态GNN的系统对比，以及统计显著性检验。

\subsection{多模型节点分类对比}

表\ref{tab:model_comparison}展示了13种模型在T+I+F多模态输入下的节点分类结果，其中GT（TransformerConv）仅参与Few-shot实验（详见5.5节），故未列入此表。

\begin{table}[htbp]
  \centering
  \caption{十三模型节点分类结果（T+I+F, Accuracy \%）}
  \label{tab:model_comparison}
  \small
  \begin{tabular}{llccc}
    \toprule
    模型 & 类型 & Accuracy (\%) & F1 (\%) & 排名 \\
    \midrule
    GCNII & 标准GNN & \textbf{91.09$\pm$0.14} & \textbf{91.09$\pm$0.16} & 1 \\
    DMGC & 多模态GNN & 90.93$\pm$0.36 & 90.90$\pm$0.34 & 2 \\
    MHGAT & 多模态GNN & 90.91$\pm$0.21 & 90.91$\pm$0.21 & 3 \\
    MMGCN & 多模态GNN & 90.91$\pm$0.18 & 90.87$\pm$0.21 & 4 \\
    GraphSAGE & 标准GNN & 90.91$\pm$0.14 & 90.89$\pm$0.14 & 5 \\
    GATv2 & 标准GNN & 90.90$\pm$0.08 & 90.88$\pm$0.05 & 6 \\
    LGMRec & 多模态GNN & 90.81$\pm$0.17 & 90.79$\pm$0.19 & 7 \\
    MGAT & 多模态GNN & 90.76$\pm$0.19 & 90.77$\pm$0.18 & 8 \\
    GCN & 标准GNN & 90.76$\pm$0.04 & 90.74$\pm$0.06 & 9 \\
    GAT & 标准GNN & 90.71$\pm$0.22 & 90.68$\pm$0.21 & 10 \\
    MLP & 标准MLP & 90.33$\pm$0.27 & 90.30$\pm$0.24 & 11 \\
    GIN & 标准GNN & 90.33$\pm$0.31 & 90.29$\pm$0.31 & 12 \\
    DGF & 多模态GNN & 90.22$\pm$0.38 & 90.13$\pm$0.41 & 13 \\
    \bottomrule
  \end{tabular}
\end{table}

\textbf{（1）GCNII以91.09\%取得全局最优。}初始残差连接与恒等映射使其在深层网络中有效传播信息，超越了所有标准GNN和多模态GNN。

\textbf{（2）多模态GNN与标准GNN性能接近。}多模态GNN组均值（90.76\%）略高于标准GNN组均值（90.63\%），但最佳标准模型GCNII（91.09\%）仍高于所有多模态GNN。排除GCNII后，DMGC（90.93\%）和MHGAT（90.91\%）超过其他标准GNN模型（GAT 90.71\%等）约0.2pp。

\textbf{（3）Top-5模型准确率差距仅0.19pp}（90.91\%-91.09\%），表明在该数据集上已接近约91\%的性能天花板。

\textbf{（4）GIN表现不佳}（90.33\%）。GIN侧重于图的结构等变特征，但在以节点特征为主导的多模态分类任务中优势不显。

\textbf{（5）链接预测方面}，由于该任务在AMMAG上高度饱和（MRR $>$ 0.97，详见5.6节），仅对部分代表性模型进行了LP评估。GCN取得最高MRR（98.50\%），但所有模型的LP差异均不足以作为模型选择的依据。

\subsection{统计显著性检验}

对12个模型各运行10次独立实验，使用成对t检验评估模型间差异的统计显著性。表\ref{tab:significance}展示各模型的显著性统计。

\begin{table}[htbp]
  \centering
  \caption{统计显著性检验结果（10次运行）}
  \label{tab:significance}
  \begin{tabular}{lrrrr}
    \toprule
    模型 & 均值 (\%) & 标准差 (\%) & 显著优于数 & 净显著胜场 \\
    \midrule
    GCNII & \textbf{91.15} & 0.30 & 8 & \textbf{+7} \\
    DMGC & 91.08 & 0.42 & 6 & +5 \\
    MHGAT & 90.99 & 0.37 & 5 & +4 \\
    MGAT & 90.97 & 0.29 & 5 & +3 \\
    GATv2 & 90.94 & 0.32 & 5 & +3 \\
    MMGCN & 90.87 & 0.31 & 4 & +2 \\
    GAT & 90.95 & 0.35 & 4 & +2 \\
    GCN & 90.77 & 0.34 & 4 & 0 \\
    LGMRec & 90.82 & 0.26 & 4 & $-$1 \\
    MLP & 90.63 & 0.35 & 2 & $-$4 \\
    GIN & 90.49 & 0.39 & 1 & $-$6 \\
    DGF & 90.33 & 0.39 & 1 & $-$8 \\
    \bottomrule
  \end{tabular}
\end{table}

\textbf{（1）GCNII具有最强的统计显著性优势}，在11对比较中8次显著优于对手，仅被GATv2显著超越1次。

\textbf{（2）Top-tier模型间差异不显著}：GCNII、DMGC、MHGAT之间均未达到$p < 0.01$的显著性水平，构成"性能无差异梯队"。

\textbf{（3）DGF统计上显著劣于几乎所有模型}，在9对比较中被显著超越。

\textbf{（4）GCN与GAT差异不显著}（$p > 0.05$），谱方法与注意力机制在该数据集上效果接近。

\section{鲁棒性实验}

本节从特征噪声、标签噪声、模态缺失和边噪声四个维度评估模型鲁棒性，全部使用T+I+F多模态输入。

\subsection{特征与标签噪声鲁棒性}

在6个模型上注入三类噪声：文本噪声（shuffle文本特征）、图像噪声（shuffle图像特征）和标签噪声（随机翻转标签），噪声比例0\%-80\%。

表\ref{tab:noise_retention}展示了80\%噪声下各模型的性能保留率。

\begin{table}[htbp]
  \centering
  \caption{80\%噪声下的性能保留率（Accuracy \% / 保留率 \%）\protect\footnotemark}
  \label{tab:noise_retention}
  \begin{tabular}{lrrrrrr}
    \toprule
    噪声类型 & GAT & GraphSAGE & DGF & MMGCN & MGAT & LGMRec \\
    \midrule
    无噪声基线 & 91.13 & 91.31 & 90.60 & 90.78 & 91.16 & 91.02 \\
    \midrule
    文本噪声 & 86.65 & 86.54 & 85.31 & 86.55 & 86.42 & 86.52 \\
    保留率 & 95.2\% & 94.9\% & 94.1\% & 95.4\% & 95.1\% & 95.1\% \\
    \midrule
    图像噪声 & 90.95 & 91.02 & 89.97 & 90.93 & 90.63 & 90.83 \\
    保留率 & \textbf{99.7\%} & \textbf{99.6\%} & 99.1\% & \textbf{99.6\%} & 99.5\% & \textbf{99.8\%} \\
    \midrule
    标签噪声 & 13.95 & 12.43 & 17.60 & 22.28 & 22.87 & 23.21 \\
    保留率 & 15.3\% & 13.6\% & 19.4\% & 24.5\% & 25.1\% & 25.5\% \\
    \bottomrule
  \end{tabular}
\end{table}

\footnotetext{鲁棒性实验使用独立的训练运行，无噪声基线与表\ref{tab:model_comparison}存在轻微差异（$\leq$0.5pp），属于正常实验波动。}

\textbf{（1）图像噪声几乎不影响性能}：80\%图像噪声下保留率$>$99\%，模型主要依赖文本和图结构，图像信息对分类是冗余信号。

\textbf{（2）文本噪声影响温和}：80\%文本噪声下保留率约95\%，模型通过图结构邻居聚合弥补了文本特征损失。

\textbf{（3）标签噪声是最大威胁}：60\%标签噪声时DGF率先崩溃（70.59\%），其余维持86-87\%；80\%时所有模型崩溃至12-23\%，呈现"悬崖效应"。

\textbf{（4）多模态GNN在特征噪声下并不比标准GNN更鲁棒}，但MMGCN对标签噪声略具韧性（80\%时22.28\% vs GAT 13.95\%）。

\subsection{模态缺失鲁棒性}

测试5种模态缺失场景下的性能退化。

\begin{table}[htbp]
  \centering
  \caption{模态缺失鲁棒性（Accuracy \%，6模型均值）}
  \label{tab:modality_missing}
  \begin{tabular}{lcc}
    \toprule
    缺失场景 & Accuracy (\%) & 相对完整损失 (pp) \\
    \midrule
    T+I+F（完整基线） & $\sim$90.70 & --- \\
    Text-only（缺失图像+公式） & 90.49 & $-$0.21 \\
    Text+Image（缺失公式） & 90.79 & +0.09 \\
    Random 50\%特征丢失 & 90.69 & $-$0.01 \\
    Image-only（缺失文本+公式） & 82.24 & $-$8.46 \\
    Formula-only（缺失文本+图像） & 80.90 & $-$9.80 \\
    \bottomrule
  \end{tabular}
\end{table}

\textbf{（1）仅保留文本（Text-only）仍可达到99.4\%的性能保留率}（90.5\%-90.9\%），几乎不影响分类能力。

\textbf{（2）缺失文本是灾难性的}：Image-only和Formula-only分别降至约82\%和81\%，损失8-10pp。

\textbf{（3）Random 50\%特征丢失被良好容忍}（约90.7\%）：图结构邻居信息提供了有效的冗余补偿。

\textbf{（4）模态贡献排序}：文本 $\gg$ 图像 $\geq$ 公式，文本信息是性能的基石。

\subsection{边噪声鲁棒性}

在6个模型上添加随机虚假边（不删除原始边），噪声比例0\%-80\%。

\begin{table}[htbp]
  \centering
  \caption{80\%边噪声下的性能保留率}
  \label{tab:edge_noise}
  \begin{tabular}{lrrrrrr}
    \toprule
    & GAT & GraphSAGE & DGF & MMGCN & MGAT & LGMRec \\
    \midrule
    无噪声基线 (\%) & 91.13 & 91.31 & 90.60 & 90.78 & 91.16 & 91.02 \\
    80\%噪声 (\%) & 75.27 & 90.68 & 55.07 & 63.89 & 79.49 & 63.71 \\
    保留率 (\%) & 82.7 & \textbf{99.4} & 60.6 & 70.3 & 87.3 & 70.0 \\
    \bottomrule
  \end{tabular}
\end{table}

\textbf{GraphSAGE对边噪声极其鲁棒}：80\%边噪声下仅损失0.63pp（91.31\%$\to$90.68\%），其采样策略天然过滤噪声边。GAT和MGAT的注意力机制可学习降低噪声边权重。DGF最敏感（保留率仅60.6\%），特征传播直接受虚假边影响。

\section{Few-shot学习实验}

Few-shot学习实验评估模型在极端少样本条件下的性能。使用5个模型（GCN、GAT、GraphSAGE、GIN、GT），7种模态组合，训练集每类仅提供$k \in \{1, 3, 5, 10, 20\}$个样本，10次独立运行。

表\ref{tab:fewshot}展示了GCN在全部7种模态组合下的Few-shot结果。其他模型（GAT、GraphSAGE、GIN、GT）呈现出与GCN相似的趋势，此处仅展示GCN的完整消融结果。

\begin{table}[htbp]
  \centering
  \caption{Few-shot学习结果（GCN, Accuracy \%，10次运行均值$\pm$标准差）}
  \label{tab:fewshot}
  \small
  \begin{tabular}{lrrrrr}
    \toprule
    模态 & k=1 & k=3 & k=5 & k=10 & k=20 \\
    \midrule
    Text & 53.8$\pm$9.1 & 67.2$\pm$4.8 & 73.0$\pm$4.2 & 78.7$\pm$2.7 & \textbf{83.4$\pm$1.0} \\
    Image & 52.6$\pm$8.9 & 68.8$\pm$2.7 & 71.4$\pm$2.0 & 73.5$\pm$1.6 & 75.4$\pm$1.3 \\
    Formula & 37.3$\pm$4.3 & 48.1$\pm$4.2 & 53.1$\pm$2.7 & 58.4$\pm$2.8 & 63.7$\pm$2.6 \\
    T+I & 55.5$\pm$8.8 & \textbf{72.7$\pm$2.2} & \textbf{75.8$\pm$2.0} & \textbf{78.9$\pm$2.1} & 81.2$\pm$1.6 \\
    T+F & 52.4$\pm$7.7 & 65.4$\pm$3.2 & 71.8$\pm$2.8 & 77.5$\pm$2.6 & 82.8$\pm$1.4 \\
    I+F & 54.0$\pm$5.9 & 68.6$\pm$3.2 & 71.2$\pm$2.1 & 74.3$\pm$2.1 & 76.4$\pm$1.5 \\
    \textbf{T+I+F} & \textbf{57.4$\pm$9.4} & 72.0$\pm$1.8 & 75.3$\pm$2.8 & 78.4$\pm$2.1 & 81.3$\pm$1.5 \\
    \bottomrule
  \end{tabular}
\end{table}

图\ref{fig:fewshot}展示了Few-shot学习曲线。

\begin{figure}[htbp]
  \centering
  \includegraphics[width=0.92\textwidth]{images/few_shot_curves.png}
  \caption{Few-shot学习曲线（7种模态组合，GCN）}
  \label{fig:fewshot}
\end{figure}

\textbf{（1）多模态组合在极端少样本（k=1）下最优}：T+I+F（57.4\%）和T+I（55.5\%）显著领先单模态，多模态增益在数据极度稀缺时最显著。随着$k$增大，T+I在$k \in \{3,5,10\}$时持续最优，Text在$k=20$时追平（83.4\%）。

\textbf{（2）Formula从最低起点实现最大相对改善}：k=1时仅37.3\%（最差），但k=20时达到63.7\%（相对提升70.8\%），公式特征需要更多样本来学习有效表示。

\textbf{（3）Text在k=20时达到最高83.4\%}，T+F紧随其后（82.8\%），两者差距仅0.6pp。结合全监督实验中T+F优于Text（90.84\% vs 89.30\%）的结论，T+F仍是最稳健的模态组合。

\textbf{（4）k=10是性价比最高的设置}：T+I达到78.9\%（仅比k=20低3pp），而训练样本仅需40个（4类$\times$10）。

\textbf{（5）Image饱和最早}：k=10后仅从73.5\%升至75.4\%（+1.9pp），远低于Text的增幅（+4.7pp）。

\section{链接预测实验}

\subsection{全链接预测与饱和分析}

如5.2.4节所示，链接预测任务在该数据集上高度饱和——所有模型的MRR $>$ 0.93，Hits@3和Hits@10均为1.0。这反映了AMMAG引用网络的强同质性：同领域论文高度互联，链接预测退化为领域识别的代理任务。

与MAGB\cite{yan2025magb}报告的MRR 0.4-0.8相比，AMMAG的饱和性能表明该任务在当前图结构上已无法区分不同方法的优劣。

\subsection{跨域链接预测}

为解决任务饱和问题，设计了跨域链接预测任务：仅使用3,838条跨领域引用边（占总引用边59,364条的6.5\%）作为正样本，负采样同样强制为跨域节点对。表\ref{tab:cross_domain_lp}展示了跨域链接预测结果。

\begin{table}[htbp]
  \centering
  \caption{跨域链接预测结果}
  \label{tab:cross_domain_lp}
  \begin{tabular}{llcc}
    \toprule
    模型 & 模态 & MRR & Hits@10 \\
    \midrule
    GCN & Text & \textbf{0.9915} & 1.0000 \\
    GCN & Image & 0.9855 & 1.0000 \\
    GCN & T+I+F & 0.9907 & 1.0000 \\
    GAT & Text & 0.9848 & 1.0000 \\
    GraphSAGE & Text & 0.9894 & 1.0000 \\
    MLP & Text & 0.9755 & 1.0000 \\
    MLP & Image & 0.9560 & 1.0000 \\
    MLP & T+I+F & 0.9777 & 1.0000 \\
    \bottomrule
  \end{tabular}
\end{table}

跨域链接预测仍然高度饱和（MRR $>$ 0.95, Hits@10 = 1.0），表明该数据集的图结构（平均度2.60、引用关系明确）使得链接预测本身过于简单。GCN+Text取得最高MRR（0.9915），与最低的MLP+Image（0.9560）仅差3.5pp。结论：AMMAG数据集不适合作为链接预测Benchmark，建议在论文中以节点分类作为主要评估指标。

\section{跨域迁移实验}

\subsection{实验设置}

在4个领域之间进行12对跨域迁移实验（每方向3个模型：GCN、GAT、GraphSAGE）。训练流程：(1) 在源域子图上预训练GNN（使用子分类标签，每域6-7类，任务难度高于4类域分类）；(2) 复制GNN编码器权重（跳过分类头）；(3) 在目标域上微调全部参数。3次独立运行。

\subsection{实验结果}

表\ref{tab:transfer}展示了12对迁移的平均增益。

\begin{table}[htbp]
  \centering
  \caption{跨域迁移平均增益（Finetune$-$Scratch, pp, 3模型均值）}
  \label{tab:transfer}
  \begin{tabular}{lrrrr}
    \toprule
    源$\to$目标 & cs.SI & cs.LO & math.LO & physics.soc-ph \\
    \midrule
    cs.SI & --- & +0.04 & $-$0.07 & $-$0.04 \\
    cs.LO & $-$0.17 & --- & $-$0.10 & +0.11 \\
    math.LO & $-$0.48 & $-$0.22 & --- & +0.42 \\
    physics.soc-ph & +0.40 & $-$0.02 & $-$0.02 & --- \\
    \bottomrule
  \end{tabular}
\end{table}

图\ref{fig:heatmap}展示了跨域迁移准确率热力图。

\begin{figure}[htbp]
  \centering
  \includegraphics[width=0.6\textwidth]{images/heatmap.png}
  \caption{跨域迁移子分类分类准确率热力图}
  \label{fig:heatmap}
\end{figure}

\subsection{结果分析}

\textbf{（1）跨域迁移增益极其有限}：平均迁移增益在$\pm$0.5pp以内，12对迁移中仅4对获得正向增益（最大+0.42pp）。预训练对目标域性能几乎没有帮助。

\textbf{（2）math.LO$\to$cs.SI迁移效果最差}（$-$0.48pp）：数学领域学习的图结构模式对计算机科学分类反而产生负迁移。

\textbf{（3）所有目标域准确率远低于全局分类}（约45-54\% vs 约91\%），验证了子分类任务（6-7类）比域分类（4类）难度大得多。

\textbf{（4）根本原因}：不同学科使用完全不同的术语体系和数学符号，子分类标签体系独立（cs.SI的子分类与math.LO的子分类无语义对应），即使同属计算机科学（cs.SI偏应用、cs.LO偏理论），域间迁移同样失败。

\textbf{（5）结论}：各领域的子分类模式差异大，GNN编码器在源域学到的表示难以迁移到不同领域，领域特定训练是必要的。

\section{编码器与训练策略分析}

本节分析编码器微调策略、编码方案选择和参数高效微调等训练策略对性能的影响。

\subsection{编码器微调策略}

在GraphSAGE+T+I+F配置下，对比四种CLIP编码器微调策略。

\begin{table}[htbp]
  \centering
  \caption{编码器微调策略对比（Accuracy \%）}
  \label{tab:encoder_tuning}
  \begin{tabular}{lc}
    \toprule
    配置 & Accuracy (\%) \\
    \midrule
    T-Frozen + V-Frozen & 90.24$\pm$0.05 \\
    T-Train + V-Frozen & 90.38$\pm$0.28 \\
    T-Frozen + V-Train & 89.94$\pm$0.05 \\
    \textbf{T-Train + V-Train} & \textbf{90.48$\pm$0.05} \\
    \bottomrule
  \end{tabular}
\end{table}

全训练（T-Train+V-Train）效果最优（90.48\%），但与全冻结仅差0.24pp。训练文本编码器增益（+0.14pp）大于视觉编码器（$-$0.30pp，反而略降）。CLIP编码器已经过充分预训练，微调的边际收益有限，推荐使用冻结策略。

\subsection{统一编码与独立编码对比}

对比统一编码器（CLIP-Large 768d，文本和图像共享语义空间）与独立编码器的跨模态检索质量。

\begin{table}[htbp]
  \centering
  \caption{统一编码与独立编码的跨模态检索对比}
  \label{tab:encoding_compare}
  \begin{tabular}{lrr}
    \toprule
    编码器 & T2I MRR (\%) & I2T MRR (\%) \\
    \midrule
    \textbf{CLIP-Large（统一）} & \textbf{11.45} & \textbf{6.56} \\
    SBERT+CLIP-Base（独立） & 0.17 & 0.20 \\
    SPECTER+CLIP-Large & 0.21 & 0.26 \\
    \bottomrule
  \end{tabular}
\end{table}

统一编码器碾压独立编码器：CLIP-Large的T2I MRR（11.45\%）是独立编码器的54-67倍。SPECTER虽然维度匹配（768d），但与CLIP不在同一语义空间，检索性能同样极低。这验证了第四章选择CLIP统一编码的合理性——跨模态对齐是多模态融合的基础。

\subsection{参数高效微调}

在GraphSAGE+T+I+F配置下对比四种参数策略。

\begin{table}[htbp]
  \centering
  \caption{参数高效微调对比（Accuracy \%）}
  \label{tab:peft}
  \begin{tabular}{lrc}
    \toprule
    PEFT策略 & Accuracy (\%) & 相对全参数 (pp) \\
    \midrule
    \textbf{None（全参数）} & \textbf{91.30$\pm$0.04} & baseline \\
    Adapter(b=64) & 89.79$\pm$0.18 & $-$1.51 \\
    LoRA(r=16) & 88.93$\pm$0.59 & $-$2.37 \\
    LoRA(r=8) & 88.12$\pm$0.24 & $-$3.18 \\
    \bottomrule
  \end{tabular}
\end{table}

全参数训练仍是最优选择。Adapter损失最小（$-$1.51pp），LoRA表现不如Adapter。GNN参数量较小（约2M），全参数训练的代价可接受，PEFT在GNN上的效果不如在LLM上显著。

\subsection{自监督基线对比}

基于GCN编码器，使用对比学习进行自监督预训练（CCA-SSG风格，不使用标签），然后在线性评估协议下测试分类性能：自监督（89.37$\pm$0.85\%）vs 有监督（91.08$\pm$0.11\%，独立运行基线），差距1.71pp。有监督训练显著优于自监督，标签信息不可替代。但自监督仍达到89.37\%，说明对比学习成功捕获了图结构和多模态特征中的有意义模式。

\section{效率分析}

表\ref{tab:efficiency}总结了各模型的训练效率。

\begin{table}[htbp]
  \centering
  \caption{模型训练效率对比\protect\footnotemark}
  \label{tab:efficiency}
  \begin{tabular}{lrrr}
    \toprule
    模型 & 参数量 & 时间/epoch (ms) & 峰值显存 (MB) \\
    \midrule
    MLP & 1.97M & 16.5 & 846 \\
    GCN & 1.97M & \textbf{13.1} & 870 \\
    GAT & 2.56M & 29.2 & 2,388 \\
    DGF & 2.03M & 21.9 & 1,663 \\
    MMGCN & 2.56M & 26.5 & 1,278 \\
    MGAT & 2.30M & 28.8 & 2,519 \\
    LGMRec & 2.04M & \textbf{13.1} & \textbf{1,008} \\
    DMGC & 2.23M & 18.8 & 1,042 \\
    MHGAT & 2.24M & 27.6 & 2,429 \\
    \bottomrule
  \end{tabular}
\end{table}

\footnotetext{GraphSAGE、GIN、GCNII、GATv2四种标准GNN模型未单独测量效率数据；其参数量与GCN相近（约2M），时间/epoch预估在13--30ms范围内。}

\textbf{（1）GCN和LGMRec训练最快}（约13ms/epoch），且LGMRec性能优秀（NC 90.81\%），是效率-性能最佳选择。

\textbf{（2）GAT/MHGAT显存消耗最大}（约2.4GB），因注意力矩阵与多头机制。

\textbf{（3）所有模型单卡可训练}：最大显存2.5GB（MGAT），远低于RTX 4090的24GB。

\textbf{（4）DMGC综合效率优异}：18.8ms/epoch + 1,042MB + 90.93\% Accuracy，适合资源受限场景。

\section{可视化分析}

\subsection{t-SNE嵌入可视化}

图\ref{fig:tsne_compare}对比了原始CLIP特征（2304维降维至2维）与GCN学习到的嵌入（256维降维至2维）的t-SNE可视化。

\begin{figure}[htbp]
  \centering
  \includegraphics[width=0.95\textwidth]{images/tsne_raw_vs_gcn.png}
  \caption{t-SNE嵌入可视化：左图为原始CLIP特征（2304d），右图为GCN嵌入（256d）}
  \label{fig:tsne_compare}
\end{figure}

原始特征空间中4个领域已有一定程度的聚类（尤其math.LO分离较好）。GCN嵌入后聚类更加紧凑，类别边界更清晰，验证了图结构聚合的有效性。cs.SI和physics.soc-ph在原始空间中重叠较多，GCN聚合后分离度改善。

\subsection{鲁棒性可视化}

鲁棒性折线图（图\ref{fig:robustness}）直观展示了三类噪声下的性能退化曲线：图像噪声曲线近乎水平（保留率$>$99\%），标签噪声曲线在60\%后急剧下降呈现"悬崖效应"，不同模型在标签噪声下的退化速率明显不同。

\begin{figure}[htbp]
  \centering
  \includegraphics[width=0.92\textwidth]{images/robustness_curves.png}
  \caption{噪声鲁棒性折线图（文本噪声、图像噪声、标签噪声）}
  \label{fig:robustness}
\end{figure}

\subsection{模态消融与Few-shot可视化}

图\ref{fig:modality_bar}展示了7种模态组合的柱状图对比，清晰可见T+F柱体最高（约91\%），Formula单模态柱体最低（约78\%），但融合后提供最大增益。

图\ref{fig:fewshot}展示了Few-shot学习曲线，Formula曲线起点最低但斜率最大（改善最快），T+I+F在k=1时优势最明显。

\subsection{跨域迁移热力图}

图\ref{fig:heatmap}展示了4$\times$4跨域迁移热力图。热力图几乎为中性色（接近0），视觉化地说明迁移增益极其有限。physics$\to$math.LO和math.LO$\to$physics有轻微正迁移，math.LO$\to$cs.SI负迁移最明显。

\section{本章小结}

本章在AMMAG数据集上进行了全面的Benchmark实验，包括多模态消融、模型有效性、鲁棒性、Few-shot学习、链接预测和跨域迁移六组核心实验，系统评估了13种模型在7种模态组合下的性能，共计产生超过1,000组实验结果。核心发现总结如下：

\textbf{（1）多模态融合在所有模型和任务上均带来提升。}节点分类提升1.05-1.66pp（T+I+F vs Text），链接预测MRR提升0.01-0.02。T+F（文本+公式）是最佳模态组合（均值90.84\%），甚至略优于三模态T+I+F（90.62\%），公式比图像提供了更多互补增益。这是AMMAG独有的发现——MAGB\cite{yan2025magb}等现有基准因缺少公式模态无法得出此结论。

\textbf{（2）不同学术领域对模态的偏好差异显著。}cs.SI从多模态融合获益最大（+2.11pp），math.LO从公式获益最大（+1.43pp），physics.soc-ph增益最小（+0.55pp）。这种差异与各领域的模态分布特征直接相关。

\textbf{（3）GCNII以91.09\%取得全局最优}，但Top-tier模型（GCNII、DMGC、MHGAT）间差异不显著（$p > 0.05$），约91\%的性能天花板需要从特征工程或数据增强层面突破。

\textbf{（4）鲁棒性实验揭示了模态贡献的层级关系。}图像噪声几乎不影响性能（保留率$>$99\%），文本是基石（Text-only保留99.4\%），缺失文本导致8-10pp性能下降。GraphSAGE对边噪声极鲁棒（保留率99.4\%），DGF最敏感（60.6\%）。

\textbf{（5）Few-shot实验发现多模态增益在数据稀缺时最显著。}T+I+F在k=1时达到57.4\%（最优），T+I在$k \in \{3,5,10\}$时持续最优，Text在k=20时追平（83.4\%）。公式特征需要更多样本释放潜力（k=1$\to$20相对提升70.8\%）。

\textbf{（6）跨域迁移增益接近零}（$\pm$0.5pp），领域间子分类模式差异大，GNN编码器难以迁移。链接预测在该数据集上饱和（MRR $>$ 0.95），不适合作为主要评估指标。

\textbf{（7）CLIP统一编码方案远优于独立编码器。}跨模态检索MRR提升67倍（11.45\% vs 0.17\%），验证了统一语义空间对多模态融合的重要性。


% 后置内容
\backmatter

% 第六章 总结与展望
%%
% 第六章 总结与展望（~2页）
%%

\begin{conclusion}

\section{工作总结}

本文针对学术多模态属性图数据集的空白，提出了AMMAG——首个面向学术论文领域的多模态属性图Benchmark数据集及其基于多智能体的自动化构建系统。本文的主要工作和贡献总结如下：

\textbf{（1）构建了首个学术多模态属性图Benchmark数据集AMMAG。}AMMAG覆盖计算机科学、数学和物理学三个学科大类，包含cs.SI（社交信息学）、cs.LO（计算逻辑）、math.LO（数理逻辑）和physics.soc-ph（社会物理）四个子领域，共计45,672篇论文。每篇论文节点包含全文文本（Markdown格式）、实验图表（PNG格式）和数学公式（LaTeX格式）三种模态属性，并通过59,364条引用边构成引用网络图结构。与现有MAG基准相比，AMMAG首次将公式模态引入MAG学习范式；与OGB\cite{hu2020ogb}和CS-TAG\cite{yan2023cstag}等学术图基准相比，AMMAG提供了原始多模态属性而非低维压缩向量表示。四个子领域在模态分布上的显著差异（如math.LO的公式密度是cs.SI的6.3倍）为模态消融实验提供了天然的实验条件。

\textbf{（2）设计了基于多智能体的自动化构建系统。}本文设计并实现了包含五个专用智能体的自动化构建系统，采用编排器-智能体-工具三层架构。五个智能体分别负责论文下载、多模态提取、数据清洗、引用图构建和标准化格式化，实现了从arXiv元数据到标准化MAG数据集的端到端自动化流程。其中，多模态提取引擎基于DocLayout-YOLO\cite{zhao2024doclayoutyolo}、UniMERNet\cite{wang2024unimernet}和PaddleOCR\cite{du2020ppocr}实现了文本、图像和公式的并行提取；引用图构建模块通过双源引用提取（PDF+Markdown，209\%提升）和三级匹配引擎（arXiv ID直接匹配、精确标题匹配、倒排索引模糊匹配）实现了自动化引用网络构建。并行化策略通过多GPU批量提取、多进程并行清洗和fork-based COW零拷贝优化实现了高效的数据处理。

\textbf{（3）进行了全面的Benchmark实验评估。}本文在多模态消融、模型有效性、鲁棒性、Few-shot学习、链接预测和跨域迁移六组任务上，系统评估了13种模型（含7种标准GNN/MLP和6种多模态GNN）在7种模态组合下的性能表现，采用CLIP-ViT-Large统一768维编码方案，共计产生超过1,000组实验结果。实验结果表明，GCNII以91.09\%取得最优节点分类准确率，但Top-tier模型间差异不显著（$p > 0.05$），约91\%的性能天花板需要从特征工程层面突破。

\textbf{（4）发现了学术领域多模态特性的重要规律。}本文首次在学术图数据集上量化分析了文本、图像和公式三种模态的贡献差异。核心发现包括：（a）T+F（文本+公式）是最佳模态组合（均值90.84\%），甚至略优于三模态T+I+F（90.62\%），公式比图像提供了更多互补增益——这是MAGB\cite{yan2025magb}等现有基准因缺少公式模态无法得出的独有发现；（b）不同学术领域对模态的偏好差异显著，cs.SI从多模态获益最大（+2.11pp）而math.LO从公式获益最大（+1.43pp）；（c）跨域迁移增益接近零（$\pm$0.5pp），领域间子分类模式差异大，领域特定训练是必要的；（d）链接预测在该数据集上饱和（MRR $>$ 0.95），不适合作为主要评估指标；（e）CLIP统一编码方案远优于独立编码器（跨模态检索MRR提升67倍）。

\section{不足与展望}

尽管本文在学术多模态属性图数据集的构建和评估方面取得了一定进展，但仍存在以下不足，每条不足对应一个未来研究方向：

\textbf{（1）数据源单一。}当前AMMAG仅覆盖arXiv单一数据源，以计算机科学、数学和物理学领域的论文为主，未涵盖生物医学、社会科学等重要学科。未来可将系统扩展至PubMed\cite{lo2020s2orc}（生物医学）、DBLP（计算机科学完整覆盖）、Semantic Scholar（跨学科）等数据源，进一步增加数据集的学科覆盖面和规模。不同数据源的论文格式和元数据标准存在差异，这需要对系统的提取和匹配模块进行适配。

\textbf{（2）引用边匹配率有待提升。}当前引用边匹配率约20.7\%，大量跨领域引用因不在候选集内而丢失。三级匹配引擎在处理OCR错误和非标准引用格式时仍存在误匹配的情况。未来可引入大语言模型（LLM）辅助引用匹配——利用LLM的语义理解能力对模糊引用进行消歧，提高匹配的准确率和召回率。此外，可以结合Semantic Scholar\cite{lo2020s2orc}和OpenAlex\cite{priem2022openalex}等外部学术知识库的引用关系数据，补充通过PDF解析未匹配到的引用边。

\textbf{（3）多模态融合策略较初步。}当前采用拼接+可学习线性投影作为多模态融合策略，实验表明T+F已接近T+I+F的性能（90.84\% vs 90.62\%），说明图像模态的互补信息有限。未来可探索更先进的融合方法：注意力融合机制（Cross-Modal Attention）能够自适应地学习不同模态的权重；多模态图Transformer可以在消息传递过程中建模模态间的交互关系；基于视觉语言模型（如LLaVA、Qwen-VL等）的统一编码器可以从根本上消除不同模态间的表示差异。

\textbf{（4）数据集规模相对较小。}AMMAG当前包含45,672篇论文，相比OGB的ogbn-arxiv（169K节点）规模较小。规模限制可能影响大规模图学习方法的评估效果。未来可通过以下途径扩展数据集规模：增加每个领域的论文数量（如从12,000篇扩展至50,000篇）、增加子领域的覆盖范围（如增加cs.AI、cs.CV等热门CS子领域）、以及引入跨arXiv的论文数据。在扩展规模的同时，需要保证多模态提取的质量——大规模数据处理对系统的并行能力和容错机制提出了更高要求。

\textbf{（5）图像编码器不适合科学图表。}虽然本文使用了CLIP统一编码方案（768维），在跨模态对齐上远优于独立编码器（检索MRR提升67倍），但CLIP视觉编码器在自然图像上预训练，其视觉特征不适用于科学图表（架构图、实验结果图、流程图等），Image-only准确率仅约81\%，远低于Text-only的约89\%。未来可训练领域适配的科学图表编码器：利用DocGenome\cite{xia2024docgenome}等大规模科学文档数据集微调视觉模型，使编码器学习科学图表的特有视觉特征；或使用多模态科学模型进行端到端的图表特征提取，将图表的视觉信息与上下文文本信息联合编码。

\textbf{（6）公式编码器仍有提升空间。}本文将LaTeX公式作为文本序列用CLIP文本编码器编码，虽已显著优于SBERT编码方案，但公式模态的单模态性能（约78\%）仍远低于文本（约89\%）。Few-shot实验发现公式特征需要更多样本释放潜力（k=1$\to$20提升26.4pp，相对提升70.8\%）。未来应引入MathBERT等数学感知编码器进行对比实验，基于公式树（AST）的编码方法可能更有效地捕获公式的结构化语义。

\end{conclusion}

% 参考文献
%%
% BIThesis 本科毕业设计论文模板 —— 使用 XeLaTeX 编译 The BIThesis Template for Undergraduate Thesis
% This file has no copyright assigned and is placed in the Public Domain.
%%

% 如无特殊需要，本页面无需更改

\begin{bibprint}

\printbibliography[heading=none]

\end{bibprint}

% 附录
%%
% BIThesis 本科毕业设计论文模板 —— 使用 XeLaTeX 编译 The BIThesis Template for Undergraduate Thesis
% This file has no copyright assigned and is placed in the Public Domain.
%%

\begin{appendices}

\section{AMMAG数据集完整实验结果}

表\ref{tab:full_nc}展示了七种模态组合的节点分类完整结果（4种基线模型，CLIP-ViT-Large 768d统一编码，5次运行均值$\pm$标准差）。

\begin{table}[htbp]
  \centering
  \caption{七种模态组合节点分类完整结果（Accuracy \%）}
  \label{tab:full_nc}
  \small
  \begin{tabular}{lccccccc}
    \toprule
    模型 & Text & Image & Formula & T+I & T+F & I+F & T+I+F \\
    \midrule
    MLP & 88.86$\pm$0.24 & 79.47$\pm$0.42 & 76.47$\pm$0.21 & 89.94$\pm$0.36 & 90.55$\pm$0.25 & 85.06$\pm$0.37 & 90.33$\pm$0.27 \\
    GCN & 89.47$\pm$0.29 & 81.25$\pm$0.28 & 78.35$\pm$0.16 & 90.37$\pm$0.16 & \textbf{90.90$\pm$0.11} & 86.29$\pm$0.37 & 90.67$\pm$0.23 \\
    GAT & 89.42$\pm$0.21 & 81.58$\pm$0.46 & 78.68$\pm$0.20 & 90.40$\pm$0.23 & 90.83$\pm$0.16 & 86.13$\pm$0.34 & 90.55$\pm$0.20 \\
    GraphSAGE & 89.43$\pm$0.22 & 81.31$\pm$0.18 & 78.43$\pm$0.12 & 90.67$\pm$0.16 & \textbf{91.06$\pm$0.22} & 86.22$\pm$0.23 & \textbf{90.91$\pm$0.14} \\
    \bottomrule
  \end{tabular}
\end{table}

表\ref{tab:full_12model}展示了13种模型在T+I+F多模态输入下的节点分类完整结果。

\begin{table}[htbp]
  \centering
  \caption{十三模型节点分类完整结果（T+I+F, Accuracy \%）}
  \label{tab:full_12model}
  \small
  \begin{tabular}{llcc}
    \toprule
    模型 & 类型 & Accuracy (\%) & F1 (\%) \\
    \midrule
    MLP & 标准MLP & 90.33$\pm$0.27 & 90.30$\pm$0.24 \\
    GCN & 标准GNN & 90.76$\pm$0.04 & 90.74$\pm$0.06 \\
    GAT & 标准GNN & 90.71$\pm$0.22 & 90.68$\pm$0.21 \\
    GraphSAGE & 标准GNN & 90.91$\pm$0.14 & 90.89 \\
    GIN & 标准GNN & 90.33$\pm$0.31 & 90.29$\pm$0.31 \\
    GCNII & 标准GNN & \textbf{91.09$\pm$0.14} & \textbf{91.09$\pm$0.16} \\
    GATv2 & 标准GNN & 90.90$\pm$0.08 & 90.88$\pm$0.05 \\
    \midrule
    DGF & 多模态GNN & 90.22$\pm$0.38 & 90.13$\pm$0.41 \\
    MMGCN & 多模态GNN & 90.91$\pm$0.18 & 90.87$\pm$0.21 \\
    MGAT & 多模态GNN & 90.76$\pm$0.19 & 90.77$\pm$0.18 \\
    LGMRec & 多模态GNN & 90.81$\pm$0.17 & 90.79$\pm$0.19 \\
    DMGC & 多模态GNN & 90.93$\pm$0.36 & 90.90$\pm$0.34 \\
    MHGAT & 多模态GNN & 90.91$\pm$0.21 & 90.91$\pm$0.21 \\
    \bottomrule
  \end{tabular}
\end{table}

\section{数据集统计信息}

表\ref{tab:full_stats}展示了AMMAG数据集的详细统计信息。

\begin{table}[htbp]
  \centering
  \caption{AMMAG数据集详细统计}
  \label{tab:full_stats}
  \begin{tabular}{lrrrrr}
    \toprule
    统计项 & cs.SI & cs.LO & math.LO & physics.soc-ph & 合计 \\
    \midrule
    论文数 & 11,929 & 11,012 & 9,555 & 13,176 & 45,672 \\
    引用边数（域内） & 4,483 & 8,174 & 12,544 & 30,325 & 55,526 \\
    平均度 & 0.75 & 1.49 & 2.63 & 4.60 & 2.43 \\
    平均图片数/篇 & 12.41 & 6.58 & 1.83 & 10.81 & 7.91 \\
    平均公式数/篇 & 180.24 & 874.70 & 1,136.50 & 188.70 & 594.79 \\
    含图片论文比例 & 98.5\% & 81.6\% & 38.2\% & 97.3\% & 81.5\% \\
    含公式论文比例 & 99.4\% & 99.8\% & 100.0\% & 99.6\% & 99.7\% \\
    \bottomrule
  \end{tabular}
\end{table}

\end{appendices}

% 致谢
%%
% BIThesis 本科毕业设计论文模板 —— 使用 XeLaTeX 编译 The BIThesis Template for Undergraduate Thesis
% This file has no copyright assigned and is placed in the Public Domain.
%%

% 致谢部分尽量不使用 \subsection 二级标题，只使用 \section 一级标题
\begin{acknowledgements}
  值此论文完成之际，首先向我的导师……
  % TODO: 请在此处填写完整的致谢内容
\end{acknowledgements}


\end{document}
